\documentclass[11pt,a4paper]{article}

\PassOptionsToPackage{hyphens,obeyspaces,spaces}{url}
\usepackage[T1]{fontenc}
\usepackage[utf8]{inputenc}
\usepackage{lmodern}
\usepackage{microtype}
\usepackage[margin=2.3cm]{geometry}
\usepackage{amsmath}
\usepackage{booktabs}
\usepackage{array}
\usepackage{tabularx}
\usepackage{multirow}
\usepackage{graphicx}
\usepackage{xcolor}
\usepackage{tikz}
\usetikzlibrary{arrows.meta,positioning,fit,calc,shapes.geometric,backgrounds}
\usepackage{pgfplots}
\usepackage{pgfplotstable}
\usepgfplotslibrary{groupplots}
\pgfplotsset{compat=1.18}
\usepackage{listings}
\usepackage[font=small,labelfont=bf]{caption}
\usepackage{enumitem}
\usepackage[section]{placeins}
\usepackage[hidelinks]{hyperref}

\definecolor{inkprimary}{HTML}{0B0B0B}
\definecolor{inksecondary}{HTML}{52514E}
\definecolor{inkmuted}{HTML}{8A8984}
\definecolor{gridline}{HTML}{E4E3DF}
\definecolor{hostfill}{HTML}{F6F5F2}
\definecolor{modehazelcast}{HTML}{2A78D6}
\definecolor{modeactivemq}{HTML}{EB6834}
\definecolor{modepersistent}{HTML}{1BAF7A}
\definecolor{modepolling}{HTML}{EDA100}
\definecolor{modeorchestrator}{HTML}{E87BA4}

\newcommand{\code}[1]{\nolinkurl{#1}}
\newcommand{\mode}[1]{\code{#1}}
\newcommand{\run}[1]{\nolinkurl{#1}}
\newcommand{\spec}[1]{\code{SPEC.md}~\S#1}
\newcommand{\repourl}{\url{https://github.com/Query-Tarantino/stage_2}}

\lstset{
  basicstyle=\ttfamily\footnotesize,
  columns=fullflexible,
  keepspaces=true,
  frame=single,
  framesep=3pt,
  rulecolor=\color{black!25},
  aboveskip=6pt,
  belowskip=4pt,
  showstringspaces=false,
  breaklines=true,
}

\pgfplotsset{
  report/.style={
    width=7.4cm,
    height=5.4cm,
    axis lines*=left,
    axis line style={inkmuted},
    tick style={inkmuted},
    major tick length=2pt,
    ymajorgrids,
    grid style={gridline},
    tick label style={font=\footnotesize, color=inksecondary},
    label style={font=\footnotesize, color=inkprimary},
    title style={font=\small, color=inkprimary},
    legend style={font=\footnotesize, draw=none, fill=none, text=inkprimary},
    legend cell align=left,
    every axis plot/.append style={line width=1.1pt},
  },
  modeline/.style={mark size=2.4pt, error bars/y dir=both, error bars/y explicit, error bars/error bar style={line width=0.7pt}, error bars/error mark options={rotate=90, mark size=2pt, line width=0.7pt}},
}

\setlist{itemsep=2pt,topsep=4pt}
\setlength{\emergencystretch}{3em}
\renewcommand{\arraystretch}{1.12}
\newcolumntype{L}[1]{>{\raggedright\arraybackslash}p{#1}}
\newcolumntype{Y}{>{\raggedright\arraybackslash}X}

\hypersetup{
  pdftitle={Search Engine Project: Stage 2. Building a Service-Oriented Architecture},
  pdfauthor={Query Tarantino: Abel Castro Villar, Hugo Gonz\'alez Fern\'andez, Javier Mart\'in Rodr\'iguez, M\'aximo Prosperi Reynoso},
}

\pgfplotstableread[col sep=comma]{
batch,replicas,mean,ci
1,1,5.52,0.05
10,1,12.52,0.21
100,1,28.51,0.85
250,1,33.28,1.84
100,2,48.38,3.58
100,3,59.75,2.21
}\indexresults
\pgfplotstableread[col sep=comma]{
clients,threads,replicas,mean_ms,ci_ms,p95_ms,rps,rps_ci
1,32,1,1.53977,0.01864,2.7328,646.26,7.79
8,32,1,2.2016999999999998,0.05982,3.8706,3618.06,99.33
32,32,1,7.28746,0.04395,14.4876,4374.39,26.49
64,32,1,14.49621,0.16188999999999998,27.228,4380.54,50.16
32,8,1,12.88946,1.25831,19.272599999999997,2493.26,254.23
32,128,1,7.55948,0.06493,15.093399999999999,4217.34,35.8
32,32,2,5.66047,0.022600000000000002,10.1446,5625.08,24.34
32,32,3,5.82144,0.04284,10.499799999999999,5471.37,39.41
}\searchresults
\pgfplotstableread[col sep=comma]{
replicas,hz,hzci,amq,amqci,amqp,amqpci,poll,pollci,orch,orchci
1,28.530,1.040,28.450,0.450,28.080,0.210,27.520,0.190,27.740,0.420
2,46.730,4.300,45.470,4.480,43.310,2.220,47.890,0.740,44.700,3.090
3,58.560,2.270,57.470,5.600,55.050,2.200,55.530,2.820,34.630,1.480
}\modesthroughput
\pgfplotstableread[col sep=comma]{
replicas,hz,hzci,amq,amqci,amqp,amqpci,poll,pollci,orch,orchci
1,51.436,1.867,52.212,0.985,50.276,0.439,54.449,0.254,55.918,0.904
2,31.942,3.113,33.493,3.355,32.401,2.015,32.062,0.592,36.203,2.138
3,25.622,0.911,26.837,2.406,24.205,0.789,28.206,1.431,31.910,2.871
}\modespublication
\pgfplotstableread[col sep=comma]{
consumers,hz,hzci,amq,amqci,amqp,amqpci
1,8872.1,583.0,9349.7,45.1,1474.7,35.2
3,8513.0,860.0,18338.0,2439.9,1456.6,10.8
8,8866.0,384.6,25904.5,3829.8,1454.0,17.8
}\brokerresults

\begin{document}

\hypersetup{pageanchor=false}
\begin{titlepage}
\centering
\vspace*{2.2cm}
{\large\bfseries Big Data - Bachelor's Degree in Data Science and Engineering\par}
\vspace{0.3cm}
{\large Universidad de Las Palmas de Gran Canaria\par}
\vspace{0.3cm}
{Academic Year 2026-2027\par}

\vspace{2.6cm}
{\fontsize{26}{32}\selectfont Search Engine Project: Stage 2\par}
\vspace{0.4cm}
{\Large Building a Service-Oriented Architecture\par}

\vspace{2.6cm}
{\large\textbf{Group Name:} Query Tarantino\par}
\vspace{0.6cm}
{\textbf{GitHub Repository:}\par}
\vspace{0.15cm}
{\repourl\par}

\vspace{2.4cm}
{\large\textbf{Team Members:}\par}
\vspace{0.4cm}
\begin{tabular}{@{}l@{\hspace{1.6cm}}l@{}}
\textbf{Abel Castro Villar} & ID: 78593645T \\
\textbf{Hugo González Fernández} & ID: 54587350R \\
\textbf{Javier Martín Rodríguez} & ID: 43386897L \\
\textbf{Máximo Prosperi Reynoso} & ID: 42252471E \\
\end{tabular}

\vfill
{November 2026\par}
\end{titlepage}
\hypersetup{pageanchor=true}

\tableofcontents
\newpage

\section*{Executive Summary}
\addcontentsline{toc}{section}{Executive Summary}

Query Tarantino turns the Stage~1 data layer into four Java~25 services that run on one host with Docker Compose.
\code{crawler} downloads Project Gutenberg books into the datalake, \code{indexer} publishes them in batches to MongoDB,
\code{query} answers keyword search with author, language and year filters, and \code{control} seeds the work and
repairs what failures leave behind. The services coordinate by choreography over Hazelcast queues; MongoDB is the only
source of truth, every state change is one conditional write, and a book becomes searchable only through a published
batch. The same code runs in four other coordination modes, and one overnight campaign measured all five on one
machine with 24 hardware threads. Table~\ref{tab:summary} answers the questions of the assignment with it.

\begin{table}[ht]
\centering
\small
\caption{The questions of the assignment and what the campaign measured, in the default mode \mode{hazelcast}.}
\label{tab:summary}
\begin{tabularx}{\textwidth}{@{}L{3.9cm}YL{1.45cm}@{}}
\toprule
Question & Answer & Section \\
\midrule
Books ingested and indexed per second & One crawler downloads 504 books/s with 8 workers and one indexer publishes 28.5 with batches of 100, so indexing limits the pipeline & \ref{sec:ingest}, \ref{sec:index} \\
Search response time under concurrent load & 7.3\,ms on average and 14.5\,ms at p95 with 32 clients; one \code{query} replica saturates at about 4{,}400 requests/s & \ref{sec:search} \\
CPU and memory as the load grows & MongoDB's CPU follows the indexers, 78--220\,\% of a core; the host never runs out of CPU, and no container's memory grows with the load & \ref{sec:resources} \\
Clients, threads and replicas & Three indexers index $\times$2.10 faster and a second \code{query} replica answers 29\,\% more; crawler replicas add nothing, because their workers wait on shared file I/O & \ref{sec:ingest}--\ref{sec:search} \\
Overhead of the coordination & A message takes 0.1--0.2\,ms and the broker uses under 17\,\% of a core; books wait for a free worker, not for the broker & \ref{sec:coordcost}, \ref{sec:modes} \\
Limits of one host, and Stage~3 & One MongoDB, one broker and one disk; Stage~3 shards the postings, spreads the datalake over the nodes and runs Hazelcast and \code{control} on several nodes & \ref{sec:limits}, \ref{sec:stage3} \\
\bottomrule
\end{tabularx}
\end{table}

All 165 recovery runs of the campaign recover with the same books and search answers in every mode and with 1, 2 and 3
replicas, and continuous integration validates every API response against \code{api/openapi.yaml} in every mode. Under
the decision rule fixed before the campaign, \mode{hazelcast} wins: it ties ActiveMQ on performance and visibility with
the least coordinator CPU. The runner-up, persistent ActiveMQ, brings lost work back without \code{control}, a property
Stage~3 adds to the Hazelcast design (Section~\ref{sec:winner}).

\newpage

\section{Introduction and Objectives}
\label{sec:introduction}

Stage~1~\cite{stage1} built the data layer of our search engine over Project Gutenberg: a datalake with the \code{batch} layout
and an inverted index in MongoDB, chosen after benchmarking three languages and several structures. Stage~2~\cite{assignment} turns that
data layer into cooperating services. Each service runs as its own process, they share explicit contracts, and an
explicit mechanism coordinates them. Everything runs on one machine, and replicas and threads are part of the design
from the start. The goal is a system whose limits we can measure on that machine and whose design we can carry to the
Stage~3 cluster.

We set six objectives and met all of them:

\begin{enumerate}
\item \textbf{Four services and one module of contracts.} \code{crawler} (Ingestion), \code{indexer} (Indexing),
  \code{query} (Search) and \code{control} (seeding and reconciliation) are written in Java~25 with Javalin and
  Jackson. They share \code{common}, which is the only place where the code names a collection, a message, a state or a
  variable. MongoDB is the source of truth, the datalake is a Docker volume, and nginx is the only entry point.
\item \textbf{An explicit, failure-proof coordination mechanism.} Work travels as messages on Hazelcast queues, MongoDB
  decides who owns each job, and \code{control} periodically repairs whatever a crash leaves behind. A book becomes
  searchable only through a published batch, and it never appears in two published batches of the same generation.
\item \textbf{Documented and validated APIs.} The front-end, operational, health and internal interfaces are kept as
  separate inventories in \code{api/openapi.yaml} (OpenAPI~3.1, version~1.11). The non-HTTP contracts are in
  \code{SPEC.md}. An automated checker validates every response against the specification in continuous integration,
  in every coordination mode.
\item \textbf{Recovery from every failure we could cause.} The failure matrix has 14 rows, each with its own script. All
  165 recovery runs of the official campaign recover, and every row gives the same books and search answers with 1, 2
  and 3 replicas and in every mode where it applies.
\item \textbf{Alternatives measured, not only argued.} The same system can coordinate in four other ways behind one
  variable: ActiveMQ with and without persistence, polling of shared state, and classical orchestration. One campaign
  on one machine compared all five against each other.
\item \textbf{Benchmarks that can be traced.} JMH microbenchmarks and nine integration experiments ran with five
  interleaved repetitions and 95\,\% confidence intervals. Every figure in this report names the run it comes from,
  and that run keeps its command, commit, configuration and raw data in the repository (Appendix~\ref{app:runs}).
\end{enumerate}

Table~\ref{tab:team} shows how the team split the work. Each member owned one piece and reviewed another member's pull
requests, and changes to the shared contracts required all four. The work advanced through seven phases, F0 to F6,
and each phase had a gate: the walking skeleton had to show a new book in \code{/search}, and F3 required the whole
failure matrix to pass. The history of \code{dev} shows the contribution of each member.

\begin{table}[ht]
\centering
\small
\caption{Pieces owned by each member. Each owner wrote the code, scripts and runs of their piece.}
\label{tab:team}
\begin{tabularx}{\textwidth}{@{}lY@{}}
\toprule
Member & Piece \\
\midrule
Hugo González Fernández (A) & \code{crawler} (Ingestion), Docker Compose and the shared \code{Dockerfile}; \code{ingest-throughput} \\
Máximo Prosperi Reynoso (B) & \code{indexer} (Indexing and batch publication); \code{index-throughput}; the four alternative coordination modes, their experiments and the campaign script \\
Abel Castro Villar (C) & \code{query} (Search), nginx, \code{api/openapi.yaml} and the \code{checks} module; the load generator and \code{search-latency} \\
Javier Martín Rodríguez (D) & \code{control} (seeding and reconciliation), the failure matrix and its 14 recovery scripts, result processing in \code{bench}; \code{coordination-cost}, \code{resource-usage}, and the official campaign on his machine \\
\bottomrule
\end{tabularx}
\end{table}
\section{System Architecture}
\label{sec:architecture}

The system has three core services and one control process, all on a single host. They never call each other in the
default mode: they exchange messages through Hazelcast and keep their data in the datalake and in MongoDB.
Figure~\ref{fig:architecture} shows the whole deployment: the front-end boundary, the internal interactions, the
replicas and threads of each process, and the shared storage.

\begin{figure}[ht]
\centering
\begin{tikzpicture}[
  font=\footnotesize,
  >={Stealth[length=4pt]},
  box/.style={draw=inksecondary, rounded corners=2pt, fill=white, align=center, inner sep=3pt},
  service/.style={box, minimum width=3.05cm, minimum height=1.25cm},
  queue/.style={draw=modehazelcast, rounded corners=1pt, fill=white, align=center, minimum width=2.25cm, minimum height=0.55cm, inner sep=2pt},
  topic/.style={draw=modehazelcast, ellipse, fill=white, align=center, inner sep=1pt, minimum width=2.4cm, minimum height=0.6cm},
  store/.style={draw=inksecondary, fill=white, rounded corners=1pt, align=left, inner sep=3pt, anchor=west},
  frontend/.style={->, draw=modehazelcast!80!black, line width=1.1pt},
  operation/.style={->, draw=inksecondary, line width=0.8pt},
  message/.style={->, draw=modehazelcast, line width=0.8pt, dashed},
  access/.style={draw=inkprimary!70, line width=0.5pt},
  lbl/.style={font=\scriptsize, text=inksecondary, inner sep=1pt},
]
\begin{scope}[on background layer]
\fill[hostfill, rounded corners=4pt] (0,0) rectangle (15.6,9.55);
\draw[inkmuted, rounded corners=4pt] (0,0) rectangle (15.6,9.55);
\end{scope}
\node[anchor=north west, font=\footnotesize\bfseries, text=inksecondary] at (0.15,9.5) {One host (Docker Compose)};

\node[box, minimum width=13.6cm, minimum height=0.6cm] (nginx) at (7.8,8.55)
  {\textbf{nginx} $\times$1 \quad 8080 front-end \quad 8081 operations \quad 8082 internal routes (mode \mode{orchestrator} only)};

\node[box] (admin) at (5.65,10.75) {Scripts and admin};
\node[box] (fe) at (12.4,10.75) {Front-end (browser)};
\draw[operation] (admin.south) -- node[lbl, right, pos=0.3] {port 8081} (admin.south |- nginx.north);
\draw[frontend] (fe.south) -- node[lbl, left, pos=0.3] {port 8080: the front-end API boundary} (fe.south |- nginx.north);

\node[service] (control) at (1.95,6.55) {\textbf{control} $\times$1 \quad :7004\\reconciliation thread\\HTTP pool};
\node[service] (crawler) at (5.65,6.55) {\textbf{crawler} $\times$N \quad :7001\\8 download threads\\HTTP pool};
\node[service] (indexer) at (9.35,6.55) {\textbf{indexer} $\times$N \quad :7002\\1 batch thread\\HTTP pool};
\node[service] (query) at (13.05,6.55) {\textbf{query} $\times$N \quad :7003\\1 snapshot thread\\HTTP pool, read only};

\draw[operation] (control.north |- nginx.south) -- node[lbl, right] {/control} (control.north);
\draw[operation] (crawler.north |- nginx.south) -- node[lbl, right] {/ingest} (crawler.north);
\draw[operation] (indexer.north |- nginx.south) -- node[lbl, right] {/index} (indexer.north);
\draw[frontend] (query.north |- nginx.south) -- node[lbl, left, align=right] {/search, /books,\\/search/filters} (query.north);

\begin{scope}[on background layer]
\fill[white, rounded corners=3pt] (2.45,3.7) rectangle (12.55,5.15);
\draw[modehazelcast, rounded corners=3pt] (2.45,3.7) rectangle (12.55,5.15);
\end{scope}
\node[anchor=south east, font=\scriptsize\bfseries, text=modehazelcast] at (12.5,3.74) {Hazelcast $\times$1 \quad :5701};
\node[queue] (b2d) at (3.8,4.55) {books-to-download};
\node[queue] (b2i) at (7.5,4.55) {books-to-index};
\node[topic] (topic) at (11.2,4.55) {published-batches};

\draw[message] ($(control.south east)+(-0.25,0)$) -- node[lbl, right, pos=0.4] {seed} (b2d.north west);
\draw[message] (b2d.north east) -- ($(crawler.south west)+(0.25,0)$);
\draw[message] ($(crawler.south east)+(-0.25,0)$) -- node[lbl, right=2pt, pos=0.4] {BookTask} (b2i.north west);
\draw[message] (b2i.north east) -- ($(indexer.south west)+(0.25,0)$);
\draw[message] ($(indexer.south east)+(-0.25,0)$) -- node[lbl, right=4pt, pos=0.45, align=left] {Batch\\Published} (topic.north west);
\draw[message] (topic.north east) -- ($(query.south west)+(0.25,0)$);

\foreach \x in {2.25, 5.95, 9.65} {
  \draw[access] (\x,5.92) -- (\x,1.25);
  \fill[inkprimary!70] (\x,1.25) circle (1.3pt);
}
\draw[access] (13.05,5.92) -- (13.05,1.25);
\fill[inkprimary!70] (13.05,1.25) circle (1.3pt);
\draw[access, ->] (1.65,5.92) -- node[lbl, left, pos=0.84] {.tmp} (1.65,2.72);
\draw[access, ->] (5.35,5.92) -- node[lbl, left, pos=0.84] {write} (5.35,2.72);
\draw[access, <-] (9.05,5.92) -- node[lbl, left, pos=0.84] {read} (9.05,2.72);
\node[store, fill=white, minimum width=10.15cm, minimum height=0.55cm] (datalake) at (0.55,2.45)
  {\textbf{datalake volume}: header and body of each book, Stage~1 \code{batch} layout};
\node[store, minimum width=14.5cm, minimum height=0.8cm] (mongo) at (0.55,0.85)
  {\textbf{MongoDB} $\times$1 \quad :27017 \quad source of truth: \code{book_states}, \code{index_state},\\
  \hphantom{\textbf{MongoDB} $\times$1 \quad :27017 \quad }\code{metadata_g<n>}, \code{postings_g<n>}, \code{batches_g<n>}};

\end{tikzpicture}
\caption{Deployment on the single Stage~2 host, in the default mode \mode{hazelcast}. Thick blue arrows are front-end
requests through nginx port 8080, gray arrows are operational requests through port 8081, dashed arrows are Hazelcast
messages and thin vertical lines are storage access; every service reads and writes MongoDB. ``$\times$N'' marks a
service that scales with \texttt{docker compose up -{}-scale}, and each box lists the threads of one replica. Not
drawn: \code{crawler} also re-enqueues retried downloads in \code{books-to-download}, and \code{control} re-enqueues
lost messages and reindexing in both queues (\spec{4.2}).}
\label{fig:architecture}
\end{figure}

\subsection{Services and shared storage}

\begin{itemize}
\item \textbf{\code{crawler} (Ingestion Service)} takes \code{BookTask} messages from \code{books-to-download}. It claims
  each book in MongoDB, downloads it from Gutenberg or from the local mirror and stores its header and body in the
  datalake with the Stage~1 \code{batch} layout. It marks the book \code{DOWNLOADED} and then enqueues it in
  \code{books-to-index}.
\item \textbf{\code{indexer} (Indexing Service)} groups \code{books-to-index} into batches of up to 100 books or 30\,s.
  It extracts the metadata of each book and counts its terms. It writes metadata and postings into the collections of
  the write generation and publishes the batch with one conditional write. Then it announces \code{BatchPublished}.
\item \textbf{\code{query} (Search Service)} keeps an immutable snapshot of the published batches of the current
  generation. It serves keyword search with author, language and year filters, TF-IDF ranking, book details and filter
  values.
\item \textbf{\code{control}} seeds the candidate books at startup and reconciles every 30\,s. It releases expired
  leases, abandons expired batches, requeues lost messages, drives index rebuilds and deletes orphan \code{.tmp} files.
  In the default mode it never drives a book step by step, so it is not an orchestrator.
\item \textbf{MongoDB} is the only source of truth. It holds \code{book_states}, one document per book, and
  \code{index_state}, which points to the generation \code{query} serves. Each generation of the index has three
  collections, \code{metadata_g<n>}, \code{postings_g<n>} and \code{batches_g<n>}.
\item \textbf{Hazelcast} carries two work queues and one topic. It is not a source of truth: a lost or repeated message
  never corrupts state, because the claim in MongoDB decides who owns each book.
\end{itemize}

\subsection{The life of a book}

Figure~\ref{fig:lifecycle} shows the states of a book and of a batch, and who writes each transition. Every arrow is
one conditional write to one MongoDB document (Section~\ref{sec:conditional-writes}). A book becomes eligible for
indexing only once its header and body are in the datalake, and it becomes searchable only when the batch that holds
it is \code{PUBLISHED} in the current generation and \code{query} has refreshed its snapshot.

\begin{figure}[ht]
\centering
\begin{tikzpicture}[
  font=\footnotesize,
  >={Stealth[length=4pt]},
  state/.style={draw=inksecondary, rounded corners=6pt, fill=white, minimum width=2.25cm, minimum height=0.65cm, font=\footnotesize\ttfamily},
  forward/.style={->, draw=inkprimary, line width=0.8pt},
  back/.style={->, draw=inksecondary, line width=0.6pt, dashed},
  fail/.style={->, draw=modeactivemq, line width=0.6pt},
  lbl/.style={font=\scriptsize, text=inksecondary, align=center, inner sep=1.5pt},
]
\node[state] (queued) at (0,0) {QUEUED};
\node[state] (downloading) at (3.5,0) {DOWNLOADING};
\node[state] (downloaded) at (7.0,0) {DOWNLOADED};
\node[state] (indexing) at (10.5,0) {INDEXING};
\node[state, draw=modepersistent, line width=0.9pt] (published) at (13.9,0) {PUBLISHED};
\node[state, draw=modeactivemq] (failed) at (7.0,-2.75) {FAILED};

\draw[forward] (queued) -- node[lbl, above=0.36cm] {claim\\(crawler)} (downloading);
\draw[forward] (downloading) -- node[lbl, above=0.36cm] {finish download\\(crawler)} (downloaded);
\draw[forward] (downloaded) -- node[lbl, above=0.36cm] {claim for a batch\\(indexer)} (indexing);
\draw[forward] (indexing) -- node[lbl, above=0.36cm] {publish batch\\(indexer)} (published);

\draw[back] (published.north) .. controls +(0,2.1) and +(0,2.1) .. node[lbl, above] {reindex into a new generation (control)} (downloaded.north);
\draw[back] (downloading.south west) .. controls +(-0.4,-0.9) and +(0.4,-0.9) .. node[lbl, below] {retryable error or\\expired lease (crawler, control)} (queued.south east);
\draw[back] (indexing.south west) .. controls +(-0.4,-0.9) and +(0.4,-0.9) .. node[lbl, below] {batch abandoned\\(control)} (downloaded.south east);
\draw[fail] (downloading.south east) -- node[lbl, right, pos=0.55] {no retries left} (failed.north west);
\draw[fail] (indexing.south) |- node[lbl, right, pos=0.3] {missing in the datalake (indexer)} (failed.east);
\draw[back] (failed.west) -| node[lbl, below, pos=0.27] {\code{POST /ingest} retries it} (queued.south);

\begin{scope}[yshift=-4.6cm]
\node[font=\footnotesize\bfseries, text=inksecondary, anchor=west] at (-1.12,0) {Batch};
\node[state] (writing) at (3.5,0) {WRITING};
\node[state, draw=modepersistent, line width=0.9pt] (bpublished) at (8.6,0) {PUBLISHED};
\node[state, draw=modeactivemq] (abandoned) at (8.6,-1.25) {ABANDONED};
\draw[forward] (writing) -- node[lbl, above=0.36cm] {publish with a live lease (indexer)} (bpublished);
\draw[fail] (writing.south) |- node[lbl, below, pos=0.75] {lease expired (control)} (abandoned.west);
\node[lbl, anchor=west, text width=4.9cm, align=left] at (9.95,-0.6) {Its books go back to \code{DOWNLOADED}, and its postings are deleted after \code{TARANTINO_GC_SECONDS}.};
\end{scope}
\end{tikzpicture}
\caption{States of a book (\code{book_states.status}) and of a batch (\code{batches_g<n>.status}), with the service
that writes each transition. Solid arrows are the normal path; dashed arrows are repairs. A \code{PUBLISHED} batch never
changes again, and an \code{ABANDONED} one is never seen by \code{query}.}
\label{fig:lifecycle}
\end{figure}

\subsection{Processes, replicas and threads}
\label{sec:threads}

A service is a logical responsibility, and each replica is a separate JVM process. Table~\ref{tab:threads} lists the
replicas and threads of each process. Request handlers never run long jobs: they only validate, write state and
enqueue, so the HTTP pool only waits on MongoDB and the broker. Within a process, threads share only an immutable
snapshot or a \code{java.util.concurrent} structure, and the code takes no lock of its own. Replicas share nothing in
memory: they meet only in MongoDB, through conditional writes, and in the broker.

\begin{table}[ht]
\centering
\small
\caption{Replicas and threads per process (\spec{2.2}). Javalin's HTTP pool grows from 8 to 250 threads unless
\code{TARANTINO_HTTP_THREADS} caps it.}
\label{tab:threads}
\begin{tabularx}{\textwidth}{@{}llX@{}}
\toprule
Process & Replicas & Threads \\
\midrule
\code{crawler} & N, 1 by default & \code{TARANTINO_PARALLEL_DOWNLOADS} download workers (8 by default), one lease heartbeat per book on a virtual thread, HTTP pool \\
\code{indexer} & N, 1 by default & One batch loop (\code{index-batches}), one lease heartbeat per open batch on a virtual thread, HTTP pool \\
\code{query} & N, 1 by default & One snapshot thread, the only one that refreshes; the topic listener only asks it to; HTTP pool that only reads the snapshot \\
\code{control} & Always 1 & One \code{reconciliation} thread for reconciliation and datalake cleanup, so they never overlap; one \code{orchestrator} thread in that mode; HTTP pool \\
\code{mongo}, \code{hazelcast}, \code{nginx} & 1 each & Their own; nginx runs one worker per CPU \\
\bottomrule
\end{tabularx}
\end{table}

The design keeps replicas and threads from interfering in four ways:
\begin{itemize}
\item \textbf{Duplicate work.} A message, an HTTP call or a poll is only a hint. The claim in MongoDB lets one worker
  take a book or a batch, and a duplicate fails its claim and is dropped.
\item \textbf{Conflicting writes.} Every state change is one conditional single-document write whose filter carries
  the expected state and the \code{claim_id}. A worker that loses its lease writes nothing else, and a write that does
  not match is never retried.
\item \textbf{Thread exhaustion.} Handlers answer in milliseconds and never wait for a job. Worker threads are a fixed
  number per replica, and a service exits instead of losing one: when a worker thread or a periodic task dies, or its
  JVM runs out of memory.
\item \textbf{Contention.} Replicas compete for the single MongoDB, the broker and the datalake volume. Batches share a
  per-batch cost among their books, postings go in \code{bulkWrite}s of 1{,}000 terms, and \code{query} reads an
  immutable snapshot without locks. Section~\ref{sec:benchmarks} measures the contention that is left.
\end{itemize}

\subsection{Coordination mechanism}
\label{sec:coordination}

The coordination is a hybrid of two mechanisms of the assignment, choreography with a broker and time-based
coordination, over one arbiter, MongoDB (Table~\ref{tab:hybrid}). The same pattern appears twice: an event for speed
and a period in case the event is lost.

\begin{table}[ht]
\centering
\small
\caption{The three parts of the default coordination and the responsibility of each.}
\label{tab:hybrid}
\begin{tabularx}{\textwidth}{@{}L{3.2cm}L{5.2cm}Y@{}}
\toprule
Mechanism & Where & Responsibility \\
\midrule
Choreography with a broker & \code{books-to-download} and \code{books-to-index} (Hazelcast \code{IQueue}); \code{published-batches} (\code{ITopic}) & Normal path: a free worker takes a book as soon as it is queued; \code{query} refreshes on each \code{BatchPublished} \\
Time-based coordination & \code{control} reconciles every 30\,s; \code{query} refreshes every 5\,s & Safety net: expired leases, abandoned batches, lost messages and lost announcements, with \code{book_states} as the persistent marker \\
Shared persistent state & MongoDB conditional writes (\spec{5.1}) & Arbiter: the claim acknowledges the message, so a duplicate fails and is dropped \\
\bottomrule
\end{tabularx}
\end{table}

The same code runs in four other modes, selected with \code{TARANTINO_COORDINATION} and one Compose override per mode
(Table~\ref{tab:modes}). Only two things change between modes: how work reaches the workers and how \code{query} learns
of a batch. The datalake, MongoDB, the conditional writes, the index, the search and the front-end API stay the same,
so the modes can be compared on equal terms (Section~\ref{sec:modes}).

\begin{table}[ht]
\centering
\small
\caption{Coordination modes (\spec{11.1}). \mode{activemq-persistent} is \mode{activemq} with persistent delivery and
acknowledgement after the job.}
\label{tab:modes}
\begin{tabularx}{\textwidth}{@{}L{2.55cm}L{3.1cm}L{3.15cm}L{2.9cm}Y@{}}
\toprule
Mode & Mechanism & Workers take work from & \code{query} learns of a batch & Lost work comes back by \\
\midrule
\mode{hazelcast} & Choreography over Hazelcast & \code{IQueue} & Topic listener and period & \code{control} requeue \\
\mode{activemq} & The same over ActiveMQ Classic 6 & JMS queues & JMS topic and period & \code{control} requeue; broker redelivery with persistence \\
\mode{polling} & Time-based, over shared state & A sorted claim in MongoDB & Period only & Nothing to requeue: waiting books are polled \\
\mode{orchestrator} & Classical orchestration by \code{control} & An internal \code{POST} into a local queue & \code{control}'s call and period & \code{control} dispatching again \\
\bottomrule
\end{tabularx}
\end{table}
\section{API Design and Contracts}
\label{sec:api}

The API has four separate inventories: front-end, operational, health and internal. Each one has a different consumer
and a different nginx port (Table~\ref{tab:inventory}). The full specification is \code{api/openapi.yaml}, in
OpenAPI~3.1~\cite{openapi}. Its \code{info.version} follows the version of \code{SPEC.md}, now 1.11, and its tags are the four
inventories. Appendix~\ref{app:endpoints} lists every operation with its responses.

\begin{table}[ht]
\centering
\small
\caption{Endpoint inventories. The front-end only sees port 8080, and only \code{query} serves it.}
\label{tab:inventory}
\begin{tabularx}{\textwidth}{@{}L{2.1cm}L{4.65cm}L{2.15cm}Y@{}}
\toprule
Inventory & Operations & Provider & Consumer \\
\midrule
Front-end, port 8080 & \code{GET /search}, \code{GET /books/{book_id}}, \code{GET /search/filters} & \code{query} replicas behind nginx & Front-end, and the \code{checks} and the load generator acting as one \\
Operational, port 8081 & \code{POST /ingest/{book_id}}, \code{GET /ingest/status/{book_id}}, \code{GET /ingest/list} & \code{crawler} & Admin and load scripts, \code{checks} \\
 & \code{POST /index/update/{book_id}}, \code{GET /index/status/{book_id}}, \code{POST /index/rebuild}, \code{GET /index/status} & \code{indexer} & Admin, benchmark scripts \\
 & \code{GET /control/status} & \code{control} & Admin, \code{bench wait} \\
Health, no nginx & \code{GET /health} on each service port & Every service & Compose healthchecks \\
Internal, port 8082 & \code{POST /internal/download}, \code{POST /internal/index}, \code{POST /internal/snapshot/refresh} & \code{crawler}, \code{indexer}, each \code{query} & \code{control}, in mode \mode{orchestrator} only \\
\bottomrule
\end{tabularx}
\end{table}

\subsection{Front-end interface}

The front-end needs three use cases: searching with optional filters and pages, showing one book, and offering the
filter values. Search is one endpoint with optional parameters, not one route per filter.

\begin{lstlisting}
GET /search?q=pride%20prejudice&author=austen&language=en&offset=0&limit=20
200 OK
{"query": "pride prejudice", "filters": {"author": "austen", "language": "en"},
 "total": 1, "offset": 0, "limit": 20,
 "results": [{"book_id": 1342, "title": "Pride and Prejudice", "author": "Jane Austen",
              "language": "en", "year": 1998}]}
\end{lstlisting}

Table~\ref{tab:search} gives the rules of each parameter. Parameters are validated in the order of the table, and the
first one that fails gives \code{400} with Problem Details naming it in \code{parameter}. A valid query that matches
nothing answers \code{200} with \code{total: 0} and an empty list. \code{total} counts every match after the filters,
while \code{results} holds only the requested page.

\begin{table}[ht]
\centering
\small
\caption{Parameters and semantics of \code{GET /search} (\spec{6.3}).}
\label{tab:search}
\begin{tabularx}{\textwidth}{@{}L{2.2cm}Y@{}}
\toprule
Parameter & Rule \\
\midrule
\code{q} & Required. Tokenized with the same tokenizer as the index, and a book must contain every word. Missing or blank: \code{400 MISSING_PARAMETER}; only stopwords: \code{200} with no results \\
\code{author} & Every word of the author, through the same tokenizer, must be among the book's \code{author_terms}: \code{austen} finds ``Jane Austen'' \\
\code{language} & ISO~639-1 code, two lowercase letters. A book with several languages matches any of them; an unknown code gives an empty \code{200} \\
\code{year} & Four digits, equal to the Gutenberg release year of the book \\
\code{offset}, \code{limit} & 0 and 20 by default; \code{offset} $\geq 0$, \code{limit} from 1 to 100 \\
Order & Descending TF-IDF, the sum of $(1 + \ln tf) \cdot \ln(N/df)$ over the query words; ties by ascending \code{book_id} \\
\bottomrule
\end{tabularx}
\end{table}

\code{GET /books/{book_id}} returns the fields of a search result plus every language and the original year. It answers
\code{404} for a book that is not in the published snapshot, even if it is being indexed. \code{GET /search/filters}
returns the languages with their book counts, the range of years and the number of searchable books, computed once per
snapshot. All three routes answer \code{503} until \code{query} holds its first snapshot. \code{query} allows CORS from
\code{TARANTINO_CORS_ORIGIN} (\code{*} by default), so a browser front-end served from another origin can call
port 8080 directly.

\subsection{Operational interface and completion semantics}
\label{sec:async}

Every request that changes state answers \code{202 Accepted} with the book's current state and a \code{Location}
header that points to it, or \code{200} if the work was already done. The caller polls that \code{Location}: an
accepted book has finished when its state is \code{PUBLISHED}, and has failed when it is \code{FAILED}, with the reason
in \code{last_error}. Handlers never download or index, so they answer in milliseconds, and a client can use a 5\,s
timeout. This replaces the synchronous \code{POST /ingest} of the assignment's example, which would hold a request
thread for a whole download.

\begin{itemize}
\item \code{POST /ingest/{book_id}} creates the book's state and enqueues it. For a book that is already queued or
  downloading, it answers \code{202} without enqueuing it again; for a \code{FAILED} book, it retries it.
\item \code{POST /index/update/{book_id}} enqueues a \code{DOWNLOADED} book. It answers \code{409 NOT_DOWNLOADED} for a
  book that is not in the datalake yet, \code{202} without a new message for one that is already \code{INDEXING}, and
  \code{200} for a \code{PUBLISHED} one.
\item \code{POST /index/rebuild} opens a rebuild of the next generation, and \code{GET /index/status} reports its
  progress in \code{books_pending}. A second request while one runs answers \code{202} with the running rebuild.
\item \code{GET /index/status} and \code{GET /control/status} report the generation, the books per state, the batches
  published, the size of the index, both queues and the last reconciliation.
\end{itemize}

Repeating any \code{POST} duplicates nothing, whether in sequence or concurrently. The state lives in MongoDB, and each
transition is a conditional write that only one caller can win.

\subsection{Internal contracts}
\label{sec:internal}

In the default mode the services make no HTTP calls to each other: Hazelcast replaces the internal calls of an
orchestrated design. The broker contract is documented like an endpoint (Table~\ref{tab:messages}). Both messages are
JSON with the same \code{ObjectMapper} as the API, so the broker never needs the Java classes.

\begin{table}[ht]
\centering
\small
\caption{Non-HTTP contract of mode \mode{hazelcast} (\spec{4.2}). Mode \mode{activemq} uses the same names as JMS
destinations.}
\label{tab:messages}
\begin{tabularx}{\textwidth}{@{}L{2.55cm}L{2.65cm}L{4.15cm}Y@{}}
\toprule
Destination & Message & Produced by & Consumed by \\
\midrule
\code{books-to-download}, queue & \code{BookTask} & \code{control} (seeding, lost messages, expired downloads), \code{POST /ingest}, \code{crawler} retries & \code{crawler} workers, one book each \\
\code{books-to-index}, queue & \code{BookTask} & \code{crawler} after a download, \code{POST /index/update}, \code{control} (lost messages, released books, reindexing) & \code{indexer}, a batch at a time with \code{drainTo} \\
\code{published-batches}, topic & \code{BatchPublished} & \code{indexer}, after publishing a batch & Every \code{query} replica \\
\bottomrule
\end{tabularx}
\end{table}

\begin{lstlisting}
BookTask        {"book_id": 1342, "enqueued_at": "2026-10-21T10:05:00Z"}
BatchPublished  {"batch_id": "6716f0c2a1b2c3d4e5f60718", "generation": 1,
                 "book_count": 100, "published_at": "2026-10-21T10:06:00Z"}
\end{lstlisting}

The acknowledgement and retry rules follow from one principle: a message is a hint, and MongoDB decides. A consumer
acknowledges a message by claiming the book with a conditional write. If the claim does not match, the message was a
duplicate or stale, and it is dropped. No ordering is assumed. An \code{IQueue} has no acknowledgement, so a message
taken by a consumer that dies before its claim is lost. \code{control} detects that case when the queue is empty while
a book has waited longer than \code{TARANTINO_REQUEUE_SECONDS}, and enqueues the book again. A lost
\code{BatchPublished} only delays visibility, because each \code{query} replica also re-reads the published batches
every 5\,s.

In mode \mode{orchestrator}, \code{control} dispatches each book with \code{POST /internal/download} or
\code{POST /internal/index}, through nginx's internal port 8082. The call only offers the task to a bounded local queue
of the replica. It answers \code{202} if the task went in, \code{400} for an invalid body, and \code{503} with
\code{Retry-After: 1} when the queue is full, and nginx then hands the call to the next replica. \code{control} stops
dispatching a stage at the first \code{503} and dispatches again any book that has not been claimed after
\code{TARANTINO_REQUEUE_SECONDS}. It announces each new batch to every \code{query} replica with
\code{POST /internal/snapshot/refresh}.

\subsection{Errors}

Every error is Problem Details (RFC~9457~\cite{rfc9457}) with \code{application/problem+json}, produced by one handler
in \code{common}. nginx formats its own errors the same way, including the \code{503} it returns when no replica
answers. Table~\ref{tab:errors} lists the codes. Every \code{503} carries \code{Retry-After: 5}, and an unexpected
exception gives \code{500} in the same format.

\begin{table}[ht]
\centering
\small
\caption{Error codes of the API.}
\label{tab:errors}
\begin{tabularx}{\textwidth}{@{}L{4.4cm}L{1.0cm}Y@{}}
\toprule
\code{code} & HTTP & When \\
\midrule
\code{MISSING_PARAMETER} & 400 & \code{q} is missing or blank \\
\code{INVALID_PARAMETER} & 400 & A malformed or out-of-range parameter or \code{book_id}, named in \code{parameter} \\
\code{BOOK_NOT_FOUND} & 404 & The book has no state or, in \code{query}, is not in the snapshot \\
\code{NOT_DOWNLOADED} & 409 & \code{POST /index/update} for a book that is not downloaded \\
\code{DEPENDENCY_UNAVAILABLE} & 503 & MongoDB or the broker do not answer, \code{query} has no snapshot yet, or nginx finds no replica \\
\bottomrule
\end{tabularx}
\end{table}

\subsection{Validation evidence}
\label{sec:checks}

The \code{checks} module runs against the live stack and validates every response, status, header and body against
the schemas of \code{api/openapi.yaml}. It also compares answers with \code{checks/expected.json}, generated from the
20 sample books. It covers everything the assignment asks for:
\begin{itemize}
\item 14 search cases: one word (13 matches), two words, author, language (\code{en}, \code{de}, \code{fr} and an
  unknown code), year, all filters combined, two pages of size one, an offset past the end, stopwords only and a query
  with no matches;
\item parameter validation in the order of Table~\ref{tab:search}, and Problem Details for unknown paths and methods;
\item book details, filter values, every operational endpoint, \code{/health} on each service port, and the
  idempotency of repeated \code{POST}s;
\item in mode \mode{orchestrator}, the internal routes, called from inside the \code{control} container.
\end{itemize}
It runs with one command after starting the stack, \code{java -jar checks/target/checks-1.0-SNAPSHOT.jar}.
Continuous integration runs it in mode \mode{hazelcast} with one \code{query} replica and again with three, and once in
each other mode. Replicas and modes must give the same answers. The README lists the commands and the expected
results.
\section{Microservice Implementation}
\label{sec:implementation}

The four services share one Maven build, one hexagonal layout and one module of contracts. Every state change is a
single conditional write to MongoDB, so a crash at any point leaves work that \code{control} can redo. This section
covers the shared parts, the indexer in depth as an example, and the other services in outline.

\subsection{Code structure and technologies}

The system is one multi-module Maven build in Java~25. It contains \code{common}, the four services, and two modules
that are never deployed: \code{bench}, with the load generator and result processing, and \code{checks}, with the API
checks. Every dependency and plugin version lives only in the parent \code{pom.xml} (Table~\ref{tab:stack}). Each
service builds one executable jar and runs on its own, outside Compose too.

\begin{table}[ht]
\centering
\small
\caption{Technologies and versions.}
\label{tab:stack}
\begin{tabularx}{\textwidth}{@{}L{2.3cm}Y@{}}
\toprule
Piece & Choice \\
\midrule
HTTP & Javalin 6.7, one port per service (7001 to 7004) \\
JSON & Jackson 2.22 with one shared \code{ObjectMapper}: snake\_case, ISO-8601 dates, unknown fields ignored. It is Javalin's default JSON mapper, which is why it replaces the Gson of the assignment's example \\
Storage & MongoDB 7.0 with the synchronous Java driver 5.13; the datalake on a Docker volume \\
Coordination & Hazelcast 5.7, one member and every service a client; ActiveMQ Classic 6.1.8 in mode \mode{activemq} \\
Packaging & One shaded jar per service and one \code{Dockerfile} on \code{eclipse-temurin:25-jre}; Docker Compose and nginx \\
Tests and benchmarks & JUnit 5.11, Testcontainers 2.0, JMH 1.37, XChart 3.8 for the benchmark charts \\
\bottomrule
\end{tabularx}
\end{table}

Each service keeps the hexagonal layout of Stage~1. \code{model} holds records and pure logic, \code{ports} the
interfaces that logic needs, \code{adapters} their MongoDB, broker, datalake and HTTP implementations, and
\code{commands} the use cases. \code{model} and \code{commands} import no Javalin, MongoDB or broker classes, so they
are tested with in-memory fakes. Each service's \code{Main} and factory wire the pieces through constructors, without
an injection framework. Every service opens its HTTP port first, so \code{/health} answers \code{503} while it waits,
and then retries MongoDB and the broker without limit. Startup order therefore does not matter.

\subsection{Shared contracts in \texorpdfstring{\code{common}}{common}}

Everything that two or more modules use has its literal name in \code{SPEC.md} and, in code, lives only in
\code{common} (Table~\ref{tab:common}). Two services therefore cannot disagree on a collection, a field, a state or a
message. \code{common} changes only in a pull request that all four members review. A write used by a single service,
such as claiming a book or publishing a batch, stays in that service's MongoDB adapter.

\begin{table}[ht]
\centering
\small
\caption{What \code{common} holds.}
\label{tab:common}
\begin{tabularx}{\textwidth}{@{}L{2.3cm}Y@{}}
\toprule
Area & Content \\
\midrule
Model & Book, batch and rebuild states, error and problem codes, and the two messages, \code{BookTask} and \code{BatchPublished} \\
Configuration & \code{TarantinoConfig} reads every \code{TARANTINO_*} variable with its default and checks its invariants, such as heartbeat $\leq$ lease / 3; an invalid value stops the service naming the variable \\
Terms & The Stage~1 \code{Tokenizer} and stopwords, used by the indexer to count terms and by \code{query} to read \code{q} and \code{author} \\
Datalake & \code{DatalakePaths}, with the \code{batch} layout and the write to \code{<final>.<uuid>.tmp} renamed with \code{ATOMIC_MOVE} \\
MongoDB & Collection and field names, the indexes, the writes ``Create the pointer'', ``Create the state'' and ``Complete a publication'', and the filter of the books that block a generation switch \\
Messaging & The \code{Broker}, \code{BookTasks} and \code{BatchAnnouncements} interfaces, with their Hazelcast and ActiveMQ implementations and clients that connect in the background and keep retrying \\
HTTP & Problem Details for every error, \code{/health}, and \code{BookIds.parse} for every \code{{book_id}} route \\
Operations & \code{TransitionLog}, one line per state change, from which the benchmarks measure; \code{CrashPoints} for the recovery tests; and \code{LeaseHeartbeat} \\
\bottomrule
\end{tabularx}
\end{table}

\subsection{Conditional writes and leases}
\label{sec:conditional-writes}

Every state change is one conditional write to a single MongoDB document. Its filter carries the state it expects and,
while a job holds the work, its \code{claim_id}. No service reads a document, decides in Java and writes it back. Two
replicas racing for the same book or batch can therefore never both win, and no transaction is needed.
\spec{5.1} lists all 23 writes. This is the one that publishes a batch:

\begin{lstlisting}
batches_g<w>.findOneAndUpdate(
  filter: {_id: B, status: "WRITING", claim_id: C,
           $expr: {$gt: ["$lease_until", "$$NOW"]}},
  update: [{$set: {status: "PUBLISHED", published_at: "$$NOW", book_ids: books}},
           {$unset: "lease_until"}],
  returnDocument: AFTER)
\end{lstlisting}

\begin{itemize}
\item \textbf{Server time.} Leases and expirations are compared with MongoDB's \code{$$NOW} inside update pipelines~\cite{mongodb},
  never with a process clock, so clock drift between containers or hosts cannot break a lease.
\item \textbf{One write, the right result.} When the next step depends on how the document ended up, the write is a
  \code{findOneAndUpdate} that returns it, never a write followed by a read. The published batch it returns becomes the
  \code{BatchPublished} message.
\item \textbf{The heartbeat.} A job that holds a lease starts a \code{LeaseHeartbeat}, which renews it every
  \code{TARANTINO_HEARTBEAT_SECONDS} (20\,s for a 60\,s lease) on its own virtual thread. Only a renewal that does not
  match loses the lease; one that throws, as when MongoDB does not answer for a moment, is retried on the next beat.
  Before each of its writes the job checks whether the lease is lost, and once it is lost it writes nothing else.
\item \textbf{Crashes.} If a process dies or MongoDB fails halfway, nothing is marked: the lease expires and
  \code{control} redoes the work. A job never retries a write that did not match; it leaves the work to whoever won.
\end{itemize}

\subsection{The indexer}
\label{sec:indexer}

The indexer turns downloaded books into metadata and postings and makes them searchable one batch at a time. Nothing a
batch writes is visible until a single write publishes it. Figure~\ref{fig:batch} shows one batch.

\textbf{Gathering a batch.} \code{BatchCollector} waits for a first message on \code{books-to-index} and then drains the
queue until it has \code{TARANTINO_INDEX_BATCH} books (100) or \code{TARANTINO_BATCH_TIMEOUT_SECONDS} (30\,s) have
passed since that first message. A book repeated in several messages is kept once, with its earliest
\code{enqueued_at}, so the benchmarks can measure its queue wait.

\textbf{Writing it.} The indexer reads the write generation, which is the target of a running rebuild or else the
current one, and ensures that generation's indexes before its first write to it. It opens the batch as \code{WRITING}
with a new \code{claim_id} and a lease, starts its heartbeat, and claims the books with one \code{updateMany} from
\code{DOWNLOADED} to \code{INDEXING}. Only the books that this write moved are indexed; another replica may have claimed
the rest. \code{HeaderParser} reads each header and the \code{Tokenizer} counts the terms of each body. A book whose
header or body is missing goes to \code{FAILED} with \code{NOT_IN_DATALAKE}. The indexer then writes \code{metadata_g<w>} with
one \code{bulkWrite} of upserts and \code{postings_g<w>} in \code{bulkWrite}s of 1{,}000 terms, publishes, marks the
books \code{PUBLISHED} and announces the batch.

\begin{figure}[ht]
\centering
\begin{tikzpicture}[
  font=\footnotesize,
  stepnode/.style={circle, draw=inkprimary, fill=white, minimum size=0.46cm, inner sep=0pt, font=\scriptsize\bfseries},
  steptext/.style={anchor=west, text width=9.0cm, inner sep=1pt, text=inkprimary},
  lease/.style={anchor=west, draw=modehazelcast, fill=modehazelcast!8, rounded corners=1pt, inner sep=2pt, font=\scriptsize, text=modehazelcast!80!black},
  crash/.style={anchor=west, draw=modeactivemq, fill=modeactivemq!8, rounded corners=1pt, inner sep=2pt, font=\scriptsize\ttfamily, text=modeactivemq!85!black},
]
\def\gap{0.6}
\draw[inkmuted, line width=0.8pt] (0,0) -- (0,{-8*\gap});
\foreach \i/\text in {
  1/{Drain \code{books-to-index}: up to 100 books or 30\,s},
  2/{Read the write generation; ensure its indexes},
  3/{Open the batch: \code{WRITING}, \code{claim_id}, lease, heartbeat},
  4/{Claim its books: \code{DOWNLOADED} to \code{INDEXING}},
  5/{Read and parse each book; a missing file fails the book},
  6/{Write \code{metadata_g<w>}: one \code{bulkWrite} of upserts},
  7/{Write \code{postings_g<w>}: \code{bulkWrite}s of 1{,}000 terms},
  8/{Publish: one conditional \code{findOneAndUpdate}},
  9/{Mark the books \code{PUBLISHED}; send \code{BatchPublished}}} {
  \node[stepnode] at (0,{-(\i-1)*\gap}) {\i};
  \node[steptext] at (0.4,{-(\i-1)*\gap}) {\text};
}
\foreach \i/\text in {5/{lease checked}, 6/{lease checked}, 7/{lease checked}, 8/{lease in the filter}} {
  \node[lease] at (9.55,{-(\i-1)*\gap}) {\text};
}
\node[crash] at (12.75,{-0.5*\gap}) {after-dequeue};
\node[crash] at (12.75,{-6*\gap}) {during-batch-write};
\node[crash] at (12.75,{-7.5*\gap}) {after-publish};
\node[crash] at (12.75,{-8.25*\gap}) {after-mark};
\node[anchor=south west, font=\scriptsize\bfseries, text=modehazelcast!80!black] at (9.55,{0.25}) {Lease checks};
\node[anchor=south west, font=\scriptsize\bfseries, text=modeactivemq!85!black] at (12.75,{0.25}) {Crash points};
\end{tikzpicture}
\caption{One batch in the indexer: nine steps, the lease checks (blue), made before each write of steps 5 to 7 and inside the
publication filter, and the four crash points of the recovery tests (red), placed where each one stops the service. \code{during-batch-write} stops after the first
postings \code{bulkWrite}, so the batch is really half-written, and \code{after-mark} stops before
\code{BatchPublished} is sent. If the lease is lost at any check, the batch stops and is never published; if no book is
left after step 5, the batch is left to expire. In both cases \code{control} abandons it and sends its books back to the
queue.}
\label{fig:batch}
\end{figure}

\textbf{One posting per term and batch.} Stage~1 kept one document per word and added books with \code{$addToSet}.
Here each batch writes its own document per term, with the batch's books in ascending order and each term frequency at
the same position. Three properties follow. A half-written batch is invisible, because \code{query} only reads
postings whose \code{batch_id} is in its snapshot of published batches. Replicas never write the same document, so
frequent words cause no contention and no document grows towards MongoDB's 16\,MB limit. And upserts make a repeated
write harmless. The cost moves to reading: a term has one document per batch that contains it, so the batch size trades
indexing throughput against the number of documents a search reads.

\textbf{Metadata.} \code{HeaderParser} keeps the title and author of Stage~1 and adds the fields of the search
filters. \code{year} is the year of the \code{Release date:} line, never the eBook number on the same line, which can
also have four digits. \code{original_year} comes from \code{Original publication:} when the header has it, which only
24 of the 2{,}800 corpus headers do. \code{languages} holds ISO~639-1 codes in header order, and \code{author_terms} the
author's words through the same \code{Tokenizer}.

\textbf{Rebuild.} Each generation of the index lives in its own \code{metadata_g<n>}, \code{postings_g<n>} and
\code{batches_g<n>}, and \code{index_state.current} says which one \code{query} serves. During a rebuild the indexer
writes into the target generation while \code{query} keeps serving the old one. \code{control} sends every book back to
be indexed and switches the pointer with one write once no book blocks the switch. Renaming three collections would
not be atomic; switching a pointer is.

\textbf{Rebuild race, closed after the campaign.} \code{POST /index/rebuild} read the current generation and then
opened the rebuild, while \code{control} switches the generation and then finishes the rebuild: two writes on each
side. A request that fell between \code{control}'s two writes answered \code{202} but built no new generation. Since
\code{SPEC.md} 1.11, ``Open a rebuild'' takes its target from the previous rebuild's target inside its own write, so
that request opens the next generation.

\subsection{Crawler, query and control}

The other three services follow the same pattern: a claim, a lease renewed by \code{LeaseHeartbeat}, conditional
writes, and a message only after the write.

\textbf{Crawler.} Each replica runs eight download workers on platform threads, and each worker takes one book at a
time from \code{books-to-download}. A worker claims the book (\code{QUEUED} to \code{DOWNLOADING}, with its
\code{claim_id}, a lease and one more attempt) and skips the download if the book's header and body are already in the
datalake.
Otherwise it reads the text from the local mirror or over HTTP, splits it into header and body at the Gutenberg markers
and writes each file to a unique \code{.tmp}, renamed with \code{ATOMIC_MOVE}, header first. ``Finish a download'' then
marks it \code{DOWNLOADED}, and only after that write is the book enqueued in \code{books-to-index}. A failed download
goes back to \code{QUEUED} while attempts are left if the error may pass (\code{NETWORK_ERROR},
\code{STORAGE_ERROR}), and straight to \code{FAILED} if it cannot (\code{NOT_FOUND}, \code{MISSING_MARKERS}). When
Gutenberg's mirror answers \code{429} or \code{503}, every download of the replica waits what \code{Retry-After} asks,
or 1, 2, 4, 8 and 16\,s.

\textbf{Query.} \code{query} keeps an immutable snapshot of the published batches of the current generation behind an
\code{AtomicReference}. One thread refreshes it at startup, on each \code{BatchPublished} and every 5\,s, comparing
batches by \code{_id}, never by publication time. A search reads the snapshot once and tokenizes \code{q}. With
filters, it first asks \code{metadata_g<n>} for the books matching \code{author_terms}, \code{languages} and
\code{year}. It then reads only the postings of the snapshot's batches, intersects the terms starting with the rarest,
scores each book with the sum of $(1 + \ln tf) \cdot \ln(N/df)$, breaks ties by \code{book_id} and reads the metadata of
the page in one query. Every request therefore sees one consistent published state, and no lock is taken on the read
path.

\textbf{Control.} At startup \code{control} creates the state of every candidate and enqueues only the books it
created, so a restart never seeds twice. Then it reconciles every \code{TARANTINO_RECONCILE_SECONDS}, in six idempotent
steps: (1) release expired downloads; (2) abandon the \code{WRITING} batches whose lease expired, in every generation;
(3) settle each \code{INDEXING} book by its batch, releasing it if the batch is abandoned and completing its publication
if the batch is published; (4) requeue the books whose message was lost, when their queue is empty, which also refills
the queues after a broker restart; (5) delete the postings of batches abandoned more than \code{TARANTINO_GC_SECONDS}
ago, on every pass, so a zombie indexer's late writes are deleted too; and (6) drive a rebuild and drop old
generations. A separate task deletes orphan \code{.tmp} files, and \code{GET /control/status} reports the generation,
the books per state, both queues and the last reconciliation.

\subsection{Coordination modes in code}

The broker sits behind three ports in \code{common}: \code{Broker} (connection and health), \code{BookTasks} (a work
queue) and \code{BatchAnnouncements} (the topic). Each service's factory picks the adapters from
\code{TARANTINO_COORDINATION}, so the commands do not know which mode runs. Each mode adds little code of its own
(non-blank lines of main code):
\begin{itemize}
\item \mode{hazelcast}: 192 lines in 7 classes, \code{IQueue} and \code{ITopic} adapters and a reconnecting client.
\item \mode{activemq}: 328 lines in 3 classes, one JMS session and producer per queue handle, with prefetch 1 on
  \code{books-to-download}. With persistence, delivery is persistent and the acknowledgement is called only when the
  job returns, so a message held by a dead consumer goes to another one.
\item \mode{polling}: 104 lines in 2 worker classes plus two sorted claims, ``Claim the next book'' and ``Claim the next
  book for a batch'', which are \code{findOneAndUpdate}s ordered by \code{enqueued_at}.
\item \mode{orchestrator}: 506 lines in 15 classes. \code{control} runs a tick every 200\,ms that dispatches the
  oldest waiting books within a window of 512 downloads and 400 indexings in flight, one \code{POST} at a time, and
  each worker replica keeps a bounded local queue behind its internal route.
\end{itemize}

\subsection{Deployment and startup}

The whole system runs on one machine with Docker Compose: \code{mongo:7.0}, one Hazelcast~5.7 member in cluster
\code{tarantino}, the four services built from one \code{Dockerfile}, and nginx. No service sets
\code{container_name}, so \texttt{docker compose up -{}-scale crawler=3 -{}-scale indexer=3 -{}-scale query=2} works,
while \code{control} always runs once. Only nginx publishes ports. It resolves the service names through Docker's DNS
(\code{127.0.0.11}, \code{valid=5s}) with the upstream in a variable, so replicas added by \texttt{-{}-scale} receive
requests without restarting it, and it retries another replica when one does not answer. Every service has a
healthcheck on \code{/health} and never restarts on its own, so the recovery tests decide when a crashed service comes
back. One shared block gives every service the same timing variables and \code{-Xmx512m -XX:+ExitOnOutOfMemoryError}.

Two lessons from building it shaped the configuration:
\begin{itemize}
\item \textbf{A timeout that is not infinite.} The Hazelcast client first used the ``infinite'' cluster connect
  timeout, $-1$, which Hazelcast stores as \code{Long.MAX_VALUE} and subtracts the elapsed time from. On a freshly
  started VM the clock stepped back a few milliseconds between two attempts, the subtraction overflowed into a negative
  wait, and the client shut down. \code{control} then died at startup whenever Hazelcast was not ready yet. The timeout
  is now ten years: in practice it never gives up, and it cannot overflow.
\item \textbf{A thread that dies quietly.} A batch keeps the term frequencies and postings of all its books in memory,
  about 0.7\,MB per book. In a pilot with batches of 500 books and a 512\,MB heap, the indexing thread threw
  \code{OutOfMemoryError} and died, but the JVM and its HTTP server did not. \code{/health} kept answering \code{UP}
  while the batch's books stayed \code{INDEXING}. Since then every JVM runs with \code{-XX:+ExitOnOutOfMemoryError}, so
  it exits and recovers like a crash, and batches stop at 250 books. The same holds for any other failure: every worker
  thread and periodic task of a service is started through \code{WorkerThreads}, which exits the service with code 1
  if anything escapes the task, instead of leaving it dead behind a healthy \code{/health}.
\end{itemize}

\subsection{Testing and continuous integration}
\label{sec:testing}

Tests follow the hexagonal layout: commands run on in-memory fakes, adapters on a real MongoDB in Testcontainers, and
every conditional write is raced by two claimants at once. \code{TwoAtOnce} releases the same call from two threads
through a barrier. For example, two indexers claiming the same books never share one, and of two publications of a
batch only one matches. Publishing and abandoning at the lease boundary never both match, and two rebuild requests open
one rebuild without an error. Commands are tested with a fake lease that is lost at a chosen step, and the tests check
that the crash points are reached in order. Routes run on a real Javalin server with fake ports, checking status
codes, \code{Location} headers, Problem Details and the exact snake\_case JSON.

On 7 November 2026, \texttt{mvn clean verify} on the delivered code (\code{8c26405}, the last commit that changes code)
ran 646 tests, all passing, in 90\,s (Table~\ref{tab:tests}). JaCoCo measures in the same build the lines and branches
each module's own tests run: 83\,\% of the lines and 83\,\% of the branches of \code{common} and the four services.
\code{checks} is low because it runs against the live stack in continuous integration, not in unit tests. GitHub Actions requires five checks before anything merges into \code{dev} or \code{main}:
the build with every test, failing if the MongoDB tests were skipped; the end-to-end API checks in every mode; the shell
and workflow linters; CodeQL; and the whole failure matrix with 1, 2 and 3 replicas on the benchmark corpus, which runs
on every pull request to \code{main} and on those to \code{dev} that change the recovery scripts.

\begin{table}[ht]
\centering
\small
\caption{Size of each module, its tests and the share of its lines they run, at \code{8c26405}. Lines include blank
lines; coverage counts executable lines, as JaCoCo does.}
\label{tab:tests}
\begin{tabular}{@{}lrrrrr@{}}
\toprule
Module & Main files & Main lines & Test lines & Tests & Line coverage \\
\midrule
\code{common} & 62 & 2{,}512 & 2{,}128 & 135 & 82\,\% \\
\code{crawler} & 31 & 1{,}311 & 2{,}245 & 93 & 82\,\% \\
\code{indexer} & 37 & 1{,}532 & 1{,}794 & 75 & 77\,\% \\
\code{query} & 20 & 1{,}063 & 1{,}046 & 47 & 89\,\% \\
\code{control} & 63 & 2{,}396 & 2{,}804 & 95 & 86\,\% \\
\code{bench} & 89 & 6{,}050 & 3{,}851 & 187 & 79\,\% \\
\code{checks} & 26 & 1{,}883 & 315 & 14 & 13\,\% \\
\midrule
Total & 328 & 16{,}747 & 14{,}183 & 646 & \\
\bottomrule
\end{tabular}
\end{table}
\section{Design Decisions}
\label{sec:decisions}

\subsection{Coordination mechanism}

We chose choreography over a broker for the normal path, MongoDB as the only arbiter of who owns each job, and a
periodic reconciliation as the safety net (Section~\ref{sec:coordination}). Three reasons drove the choice:
\begin{itemize}
\item \textbf{No process in the data path.} A free worker takes a book as soon as it is queued, and replicas scale by
  competing for the same queue. No central process calls every service for every book, so none can become the
  bottleneck or the single point of failure of the pipeline.
\item \textbf{Exactly one winner without transactions.} Messages can be lost or repeated, so they cannot decide
  anything. A conditional single-document write gives exactly one winner per transition, without locks or
  transactions, and it is valid across machines because it compares against MongoDB's clock.
\item \textbf{Recovery from persistent state.} Every repair works from \code{book_states}, so a lost message, an expired
  lease or a lost announcement is found and fixed by reading MongoDB. Hazelcast can restart with empty queues.
\end{itemize}

The design has costs, and we accepted them knowingly. It has more moving parts than an orchestrator: leases, the
heartbeat and a six-step reconciliation. It runs two systems, Hazelcast and MongoDB, where polling MongoDB alone would
work. Batching adds wait, because a book waits for its batch to fill or for 30\,s. And the safety net is slow on
purpose: \code{TARANTINO_REQUEUE_SECONDS} is 120\,s, and it must stay above the batch timeout so that \code{control}
never requeues books an indexer is still holding.

\subsection{The five modes measured, and the winner}
\label{sec:winner}

The assignment allows any coordination mechanism, so we also built the four alternatives of Table~\ref{tab:modes} and
measured all five in one campaign (Section~\ref{sec:modes}). The decision rule was written on 5~November, before the
campaign ran:
\begin{enumerate}
\item Correctness is a gate: every applicable row of the failure matrix must recover with the same answers.
\item A difference only counts when \code{pairs.csv} marks it as significant, that is, when the 95\,\% confidence
  interval of the difference between two modes excludes zero.
\item The criteria are then applied in this order: pipeline performance (throughput and time to publication),
  visibility, cost (the CPU of the coordinator and of MongoDB, the containers and the code of the mode), robustness
  (lost work, a broker restart, \code{control} down), and fit for Stage~3.
\item The default stays \mode{hazelcast} in this delivery, and a mode the data favours is proposed for Stage~3.
\end{enumerate}

Table~\ref{tab:decision} applies the rule. Every mode passes the gate: all 165 recovery runs recovered, and
\code{recovery-compare} found the same books and answers for each row in every mode.

\begin{table}[ht]
\centering
\scriptsize
\setlength{\tabcolsep}{3.5pt}
\caption{Decision matrix of the five coordination modes, in the order of the decision rule. Pipeline, visibility and
CPU figures are means over the 5 repetitions of \run{20261106-012054-coordination-variants-c1376c7} of each
repetition's value, as \code{pairs.csv} holds them: a p50 here is the mean of the five repetitions' medians, not the
median of their pooled values that \code{summary.csv} gives, so the two differ slightly. Recovery figures come from
the 165 recovery runs and \run{20261106-073911-recovery-delay-c1376c7}. ``N'' is the number of \code{crawler} and
\code{indexer} replicas together, and ``*'' marks a significant difference against \mode{hazelcast}.}
\label{tab:decision}
\begin{tabularx}{\textwidth}{@{}L{2.45cm}YYYYY@{}}
\toprule
Criterion & \mode{hazelcast} & \mode{activemq} & \mode{activemq-persistent} & \mode{polling} & \mode{orchestrator} \\
\midrule
Gate: rows recovered & 14 of 14 & 14 of 14 & 14 of 14 & 11 of 11 & 13 of 13 \\
\addlinespace
Throughput, books/s, N = 1 / 2 / 3 & 28.5 / 46.7 / 58.6 & 28.4 / 45.5 / 57.5 & 28.1 / 43.3 / 55.0* & 27.5 / 47.9 / 55.5 & 27.7 / 44.7 / 34.6* \\
Time to publication p50, s, N = 1 / 2 / 3 & 50.0 / 30.6 / 24.6 & 50.9 / 32.5 / 25.7 & 49.1 / 31.6 / 23.1* & 53.2* / 30.8 / 26.8* & 54.7* / 35.3* / 30.1* \\
\addlinespace
Visibility p50, ms & 38--45 & 39--48 & 38--48 & 2{,}279--2{,}712* & 168--211* \\
\addlinespace
Coordinator CPU, \% of a core, N = 1 / 3 & 3.6 / 8.3 & 5.4* / 12.8* & 7.8* / 16.6* & 3.6 / 7.0* & 10.8* / 12.7* \\
MongoDB CPU, N = 3 & 224\,\% & 219\,\%* & 219\,\% & 247\,\%* & 140\,\%*, doing less work \\
Coordination containers & 1, Hazelcast & 1, ActiveMQ & 1, ActiveMQ with a disk store & None & None \\
Main code of the mode & 192 lines & 328 lines & 328 lines & 104 lines & 506 lines \\
\addlinespace
Lost work (rows 04, 05) & \code{control} requeues after 11.7--11.8\,s & \code{control} requeues after 11.7--11.8\,s & Broker redelivers; reclaimed in 3.8--5.0\,s & Nothing to lose & \code{control} redispatches after 9.7\,s \\
Broker restart (row 12) & Queues lost; download messages requeued after 12\,s, index messages once \code{books-to-index} empties, 82--207\,s & Queues lost; download messages requeued after 11.7\,s & Messages survive; reclaimed in 6.0\,s & No broker & No broker \\
\code{control} down (row 11) & Pipeline goes on & Pipeline goes on & Pipeline goes on & Pipeline goes on & Pipeline stops \\
\addlinespace
Stage~3 fit (argued, not measured) & Native cluster with partition backups, in memory & Network of brokers or a shared store & A shared or replicated disk store; capped at 1{,}460 messages/s & MongoDB load grows with every worker & One dispatcher: a bottleneck and a single point of failure \\
\midrule
Verdict & \textbf{Winner} & Ties on performance, costs more & \textbf{Runner-up}: best recovery & Out on visibility & Out on performance \\
\bottomrule
\end{tabularx}
\end{table}

\textbf{Performance.} \mode{orchestrator} loses: 41\,\% less throughput with 3 replicas, because one thread dispatches
every book one \code{POST} at a time, and a median time to publication 4.7--5.5\,s longer at every N. \mode{polling}
keeps the throughput, but its median time to publication is 2.2--3.2\,s longer with 1 and 3 replicas.
\mode{hazelcast} and \mode{activemq} never differ significantly. \mode{hazelcast} and \mode{activemq-persistent} only
differ with 3 replicas, in opposite directions: \mode{hazelcast} publishes 6\,\% more books per second, and
\mode{activemq-persistent} has a 1.6\,s (6\,\%) lower median time to publication, 23.05 against 24.62\,s in
\code{pairs.csv}, which Table~\ref{tab:decision} rounds to 23.1 and 24.6. Its median is lower because persistent sends
feed the crawlers more slowly than they take books. A book waits 1\,ms in \code{books-to-download} instead of 0.8\,s, and slightly less in
\code{books-to-index}, so part of the gain is a later start of each book's clock, not faster work. The same throttling
costs the 6\,\% of throughput. Performance therefore leaves three modes in a tie.

\textbf{Visibility and cost.} \mode{hazelcast}, \mode{activemq} and \mode{activemq-persistent} make a batch searchable
within 38--48\,ms, so visibility does not separate them. Cost does: \mode{hazelcast}'s coordinator uses significantly
less CPU than both ActiveMQ modes at every N, 35\,\% less than \mode{activemq} and 50\,\% less than
\mode{activemq-persistent} with 3 replicas. It also needs no disk store and has the least code of the three. The rule
stops here: \textbf{\mode{hazelcast} wins}, and it remains the mode of Stage~2.

\textbf{What the rule leaves to Stage~3.} Robustness, the criterion after cost, favours \mode{activemq-persistent}.
It is the only mode whose lost work comes back without \code{control}: the broker redelivers what a dead consumer did
not acknowledge, and the book is claimed again 3.8--6.0\,s after the failure, against 11.7\,s for a requeue by
\code{control}. If robustness came before cost, \mode{activemq-persistent} would win. That is why its property, not its
broker, goes into the Stage~3 plan (Section~\ref{sec:stage3}). \code{control} stays necessary in every mode, because it
also releases expired leases and drives rebuilds, and persistence caps the broker at about 1{,}460 messages/s. At the
500--670 books/s that ingestion reaches, with two messages per book, that cap is close.

\subsection{Single-machine deployment}

Everything runs on one host with Docker Compose, as the assignment requires. Addresses are Compose service names, the
datalake is a local volume, and there is one MongoDB server and one broker. Ports, addresses, data paths, replica
counts, thread counts and every timing are environment variables with working defaults, validated at startup.
Replicas are added with \texttt{-{}-scale} and need no restart, because nginx resolves service names on every request
through Docker's DNS. Only nginx publishes ports, and the front-end and operational ports are separate, so a front-end
can never reach an administrative route.

\subsection{API design}

\begin{itemize}
\item \textbf{One search endpoint} with optional filters and pagination, as the assignment asks, plus book details and
  filter values for the front-end use cases.
\item \textbf{Asynchronous state changes}: \code{202} with a \code{Location} to poll, or \code{200} if already done.
  Handlers stay in the millisecond range, cannot exhaust the HTTP pool, and repeated or concurrent requests duplicate
  nothing.
\item \textbf{Consistent reads}: every search reads one immutable snapshot of published batches, so a request never
  mixes two states of the index, and a rebuild switches the whole snapshot at once.
\item \textbf{Problem Details for every error}, including nginx's, with a stable \code{code} and the failing
  \code{parameter}, so a client handles errors without parsing text.
\end{itemize}

\subsection{Storage}

The datalake keeps the Stage~1 \code{batch} layout, which won the Stage~1 comparison: $O(1)$ store and lookup with one
directory per 1{,}000 ids. The inverted index stays in MongoDB, also chosen in Stage~1, but with one postings document
per term and batch instead of one per term (Section~\ref{sec:indexer}). Metadata moves from SQLite to MongoDB: SQLite
allows one writer at a time, which several indexer replicas cannot share, and MongoDB avoids running a third database
for it. Book states also live in MongoDB, because conditional single-document writes are exactly what claims and leases
need. Each generation of the index has its own collections, so a rebuild never disturbs what \code{query} serves.

\subsection{Concurrency model and configuration}

Replicas coordinate only through conditional writes and the broker, and threads inside a process share only immutable
snapshots and \code{java.util.concurrent} structures, so no lock is ever taken. Table~\ref{tab:config} gives the
configuration we keep and the measurement behind each setting.

\begin{table}[ht]
\centering
\small
\caption{Configuration decisions and the evidence behind them (Section~\ref{sec:benchmarks}).}
\label{tab:config}
\begin{tabularx}{\textwidth}{@{}L{3.0cm}L{3.1cm}Y@{}}
\toprule
Setting & Kept & Evidence \\
\midrule
Books per batch & 100; 250 with a larger heap & 250 gives 17\,\% more books/s than 100, but the indexer container peaks at 745--808\,MiB with \code{-Xmx512m}; 500 ran out of memory \\
\code{indexer} replicas & As many as MongoDB allows & Indexing limits the pipeline; 2 and 3 replicas give $\times$1.70 and $\times$2.10, while MongoDB's CPU grows faster than the indexers' \\
\code{crawler} replicas and workers & 1 replica, 8 to 16 workers & Replicas add nothing (521--548 against 504 books/s); 16 workers add 24\,\% over 8; one crawler already downloads about 18 times faster than one indexer publishes \\
\code{query} replicas & 2 & The second replica adds 29\,\% requests/s and cuts latency by 22\,\%; the third adds nothing \\
HTTP threads & Javalin's default, 8 to 250 & 128 threads add nothing over 32 and lose 4\,\% of the requests/s; 8 threads with 32 clients raise latency by 77\,\% \\
Batch announcements & Topic listener plus a 5\,s refresh & A batch is visible in 44\,ms on average with the listener and 2.6\,s without it, at no measurable CPU cost \\
\bottomrule
\end{tabularx}
\end{table}

\subsection{Smaller decisions}

Table~\ref{tab:small-decisions} lists the smaller decisions that shaped the code, each with the alternative it
rejected.

\begin{table}[!ht]
\centering
\small
\caption{Smaller decisions, with the alternative each one rejected.}
\label{tab:small-decisions}
\begin{tabularx}{\textwidth}{@{}L{3.6cm}YL{3.4cm}@{}}
\toprule
Decision & Why & Rejected \\
\midrule
Java~25 & Allowed by the course; virtual threads carry the lease heartbeats cheaply, and records model messages and snapshots & Java~17 of the example \\
Jackson & Javalin's default mapper; one \code{ObjectMapper} for the API, MongoDB documents and broker messages & Gson of the example \\
Generation pointer for rebuilds & Switching it is one atomic write & Renaming three collections, which is not atomic \\
One conditional write per transition, on \code{$$NOW} & Single-document atomicity needs no transactions, and MongoDB's clock is the only clock & Read, decide in Java, write \\
A lease is lost only by a renewal that does not match & A moment without MongoDB should not cost a batch & Losing the lease on any exception \\
A service exits when a worker thread dies or its JVM runs out of memory & A dead thread behind a healthy \code{/health} is never recovered & Catching \code{Error} in the worker and going on \\
Ten-year Hazelcast connect timeout & The infinite one overflows when the clock steps back & $-1$ \\
JMH with three forks and the protocol in annotations & The error covers differences between JVMs, and every way of launching keeps the protocol & One fork, protocol in a runner \\
\bottomrule
\end{tabularx}
\end{table}
\section{Benchmarks and Results}
\label{sec:benchmarks}

\subsection{Setup and methodology}

Every figure in this section comes from one official campaign on one machine (Table~\ref{tab:host}). The script
\code{scripts/bench/all.sh} ran every experiment and the whole failure matrix in the five modes, in 187 steps and
9\,h~45\,min, from a clean commit of \code{dev}. Only one step failed: \code{broker-latency}, because Git Bash rewrote
the folder it passed to the container. Once the script passed a relative path, that step ran again alone, from a commit
that only changed scripts. Appendix~\ref{app:runs} gives the run, command and raw data behind every table and figure.

\begin{table}[ht]
\centering
\small
\caption{Host and fixed conditions of every official run, as each \code{manifest.json} and \code{corpus.json} record
them.}
\label{tab:host}
\begin{tabularx}{\textwidth}{@{}L{2.6cm}Y@{}}
\toprule
Host & Intel Core Ultra 9 275HX, 24 hardware threads (8 P-cores and 16 E-cores), 31.4\,GiB, 1\,TB disk, Windows~11 \\
Containers & Docker Desktop 29.9.2 with 24 CPUs and 15.3\,GiB; every service JVM with \code{-Xmx512m} \\
Corpus & 2{,}800 books and 755{,}974 distinct terms from the local mirror, checked against \code{workload/corpus.sha256} before every run, so the network is never measured \\
Campaign & \run{20261105-215348-campaign-c1376c7}, from \code{dev} at \code{c1376c7}, overnight on 5--6 November 2026 \\
\bottomrule
\end{tabularx}
\end{table}

\begin{itemize}
\item \textbf{One factor at a time.} Each experiment varies one factor around a baseline, and its script prints the
  hypotheses to the run's \code{log.txt} before measuring.
\item \textbf{Five interleaved repetitions.} Repetition 1 of every configuration runs first, then repetition 2, and so
  on, so any drift over the night, such as heat, spreads over every configuration instead of falling on one.
\item \textbf{Statistics.} Tables give the mean $\pm$ the half-width of its 95\,\% confidence interval over the
  repetitions, with Student's $t$ for $n-1$ degrees of freedom. Percentiles pool the values of every repetition: books,
  \code{docker stats} samples or requests.
\item \textbf{Three levels of concurrency, kept apart.} Client concurrency is the closed-loop clients of the load
  generator; threads are the workers or HTTP threads inside one process; replicas are processes. Each table states the
  three, and each trial records them in \code{manifest.json}.
\end{itemize}

To fit the campaign into one night, it reused two runs from 3 and 4 November on the same machine, because
\code{git diff} shows that the code they measure did not change. These are \code{search-latency} and the indexer's JMH.
The coordination cost comes from the \mode{hazelcast} trials of \code{coordination-variants}, which follow the same
procedure. Rows 09 to 12 of the failure matrix ran with one replica only outside \mode{hazelcast}, because their time is
mostly processing the whole corpus, which \code{coordination-variants} already measures for every mode and replica
count.

\subsection{Microbenchmarks}

JMH~\cite{jmh} runs in each service's \code{src/test/java} with 5 warmup and 10 measurement iterations, three forks and a fixed
heap, so its error also covers differences between JVMs. Each bound in Table~\ref{tab:jmh} was stated before measuring
and is checked with at least three sizes.

\begin{table}[ht]
\centering
\small
\caption{JMH results, average time per operation $\pm$ JMH's 99.9\,\% error. Indexer:
\run{20261103-204454-index-throughput-6925d4e}; crawler: \run{20261105-215351-ingest-throughput-c1376c7}; query:
\run{20261104-140900-search-latency-8e1c1e8}.}
\label{tab:jmh}
\begin{tabularx}{\textwidth}{@{}L{3.0cm}L{3.0cm}rL{2.15cm}Y@{}}
\toprule
Benchmark & Input & Time & Per unit & Bound \\
\midrule
\multirow{3}{3.0cm}{\code{Tokenizer}, terms of a body} & Book 1080, 20\,KB body & 0.206 $\pm$ 0.004\,ms & 10.4\,ms/MB & \multirow{3}{=}{Linear in the body: holds} \\
 & Book 84, 429\,KB & 4.83 $\pm$ 0.01\,ms & 11.3\,ms/MB & \\
 & Book 2600, 3.34\,MB & 36.7 $\pm$ 0.1\,ms & 11.0\,ms/MB & \\
\addlinespace
\multirow{3}{3.0cm}{\code{HeaderParser}} & Book 1080 & 6.9 $\pm$ 0.9\,\textmu s & -- & \multirow{3}{=}{Constant: holds} \\
 & Book 84 & 7.54 $\pm$ 0.09\,\textmu s & -- & \\
 & Book 2600 & 7.09 $\pm$ 0.03\,\textmu s & -- & \\
\addlinespace
\multirow{4}{3.0cm}{Writing a batch to MongoDB} & 1 book & 19.8 $\pm$ 0.7\,ms & 19.8\,ms/book & \multirow{4}{=}{Cost per book falls with the batch: holds} \\
 & 10 books & 338 $\pm$ 3\,ms & 33.8\,ms/book & \\
 & 100 books & 2{,}806 $\pm$ 33\,ms & 28.1\,ms/book & \\
 & 1{,}000 books & 12{,}535 $\pm$ 175\,ms & 12.5\,ms/book & \\
\addlinespace
\multirow{3}{3.0cm}{Crawler: split header and body} & Book 1080, 40\,KB & 102 $\pm$ 10\,\textmu s & -- & \multirow{3}{=}{Linear in the text} \\
 & Book 84, 450\,KB & 1.62 $\pm$ 0.01\,ms & -- & \\
 & Book 2600, 3.4\,MB & 12.3 $\pm$ 0.4\,ms & -- & \\
\addlinespace
\multirow{3}{3.0cm}{Query over 1{,}000 books} & TF-IDF ranking, 3 words & 60.2 $\pm$ 0.4\,\textmu s & -- & \multirow{3}{=}{About 1\,\% of a request end to end} \\
 & Search without filters & 44.9 $\pm$ 0.2\,\textmu s & -- & \\
 & Search with filters & 12.3 $\pm$ 0.1\,\textmu s & -- & \\
\bottomrule
\end{tabularx}
\end{table}

The tokenizer costs 10.4--11.3\,ms per MB across a 168-fold range of sizes, so it is linear, and the header parser
takes about 7\,\textmu s whatever the book. Writing a batch costs less per book from 10 to 1{,}000 books, because books
share terms and each term gets one postings document per batch. The single book is cheaper only because it is book 1,
an 11\,KB body. Splitting a book in the crawler takes at most 12\,ms, and ranking takes 60\,\textmu s. None of these
operations is what limits the pipeline, as the next sections show.

\subsection{Ingestion throughput}
\label{sec:ingest}

\begin{table}[ht]
\centering
\small
\caption{Books downloaded per second by \code{crawler} replicas and download workers per replica, over the whole window
(first claim to last download) and the steady window (10\,\% to 90\,\% of the books).
\run{20261105-215351-ingest-throughput-c1376c7}.}
\label{tab:ingest}
\begin{tabular}{@{}rrrr@{}}
\toprule
Replicas & Workers per replica & Whole window (books/s) & Steady 10--90\,\% (books/s) \\
\midrule
1 & 4 & 370 $\pm$ 51 & 393 $\pm$ 51 \\
1 & 8 & 504 $\pm$ 135 & 560 $\pm$ 132 \\
1 & 16 & 623 $\pm$ 27 & 671 $\pm$ 44 \\
2 & 8 & 548 $\pm$ 73 & 603 $\pm$ 84 \\
3 & 8 & 521 $\pm$ 17 & 572 $\pm$ 17 \\
\bottomrule
\end{tabular}
\end{table}

\begin{itemize}
\item \textbf{Workers help with a shrinking gain} (Table~\ref{tab:ingest}): +36\,\% from 4 to 8 and +24\,\% from 8 to
  16. The curve bends but does not flatten by 16 workers, as the hypothesis expected.
\item \textbf{Replicas do not help}, against the hypothesis: 2 and 3 replicas of 8 workers download as many books per
  second as one, within their intervals.
\item \textbf{The workers wait on the file I/O they share} (Table~\ref{tab:ingest-workers}). In every trial the
  workers are busy for 94--100\,\% of the window. Between two books, the enqueue, the poll of \code{books-to-download}
  and the claim in MongoDB take 1.5--2.9\,ms from 4 to 24 threads, so neither the queue nor MongoDB holds ingestion
  back. The work of a book grows instead, from 9.4\,ms with 4 threads to 23.6\,ms with 16 in one replica and 40.6\,ms
  with 3 replicas of 8, while the CPU it costs within one replica stays at 6.7--7.5\,ms: the threads mostly wait. The
  host has CPU to spare, 956--1{,}281\,\% of its 2{,}400\,\% in use (Table~\ref{tab:resources}). What that work shares is
  file I/O: reading the book from the mirror, a folder of the Windows host mounted into Docker Desktop's VM, and writing
  its header and body to the datalake volume. These runs cannot tell the read from the write; a run with the mirror
  inside a Docker volume would.
\item \textbf{Replicas cost CPU without adding books.} With 2 and 3 replicas a book costs 10.2 and 12.0\,ms of CPU,
  because each JVM warms up its own code within a window of about 5\,s. That is why 16 workers in one replica download
  more than 3 replicas of 8.
\end{itemize}

\begin{table}[ht]
\centering
\small
\caption{How the download workers of Table~\ref{tab:ingest} spend their time: the work of a book, from its claim to its
download by the same thread, the gap until that thread's next claim, and the crawlers' CPU during the ingestion, in
total and per book (the total divided by the books per second). Mean $\pm$ 95\,\% interval of 5 repetitions.
\code{workers.csv} and \code{summary.csv} of \run{20261105-215351-ingest-throughput-c1376c7}.}
\label{tab:ingest-workers}
\begin{tabular}{@{}rrrrrr@{}}
\toprule
Replicas & Workers & Work per book (ms) & Gap (ms) & Crawlers' CPU (\%) & CPU per book (ms) \\
\midrule
1 & 4 & 9.4 $\pm$ 1.3 & 1.5 $\pm$ 0.3 & 247 $\pm$ 14 & 6.7 \\
1 & 8 & 14.7 $\pm$ 4.6 & 1.8 $\pm$ 0.4 & 354 $\pm$ 71 & 7.0 \\
1 & 16 & 23.6 $\pm$ 1.1 & 2.1 $\pm$ 0.1 & 469 $\pm$ 24 & 7.5 \\
2 & 8 & 26.2 $\pm$ 4.2 & 2.6 $\pm$ 0.2 & 558 $\pm$ 86 & 10.2 \\
3 & 8 & 40.6 $\pm$ 1.6 & 2.9 $\pm$ 0.1 & 625 $\pm$ 51 & 12.0 \\
\bottomrule
\end{tabular}
\end{table}

\subsection{Indexing throughput and latency}
\label{sec:index}

\begin{figure}[ht]
\centering
\begin{tikzpicture}
\begin{groupplot}[
  group style={group size=2 by 1, horizontal sep=1.7cm},
  report,
  width=7.6cm, height=5.2cm,
  ymin=0,
  ylabel={Books per second},
  error bars/y dir=both, error bars/y explicit,
]
\nextgroupplot[
  xmode=log, log basis x=10,
  xtick={1,10,100,250}, xticklabels={1,10,100,250},
  xmin=0.8, xmax=420,
  xlabel={Books per batch, 1 replica},
  ymax=40,
  title={Batch size},
]
\addplot[color=modehazelcast, mark=*, mark size=2.4pt, error bars/error bar style={line width=0.7pt},
  nodes near coords, point meta=y, every node near coord/.style={font=\scriptsize, color=inkprimary, anchor=north west, /pgf/number format/fixed, /pgf/number format/fixed zerofill, /pgf/number format/precision=1}]
  table[x=batch, y=mean, y error=ci, restrict expr to domain={\thisrow{replicas}}{1:1}] {\indexresults};
\nextgroupplot[
  xtick={1,2,3},
  xmin=0.7, xmax=3.3,
  xlabel={\code{indexer} replicas, batches of 100},
  ymax=90,
  title={Replicas},
  clip=false,
]
\addplot[color=inkmuted, dashed, line width=0.8pt, mark=none, forget plot] coordinates {(1,28.51) (2,57.02) (3,85.53)};
\node[font=\scriptsize, color=inksecondary, anchor=south east] at (axis cs:1.95,58) {linear scaling};
\addplot[color=modehazelcast, mark=*, mark size=2.4pt, error bars/error bar style={line width=0.7pt}]
  table[x=replicas, y=mean, y error=ci, restrict expr to domain={\thisrow{batch}}{100:100}] {\indexresults};
\node[font=\scriptsize, color=inkprimary, anchor=north west] at (axis cs:1.03,27) {28.5};
\node[font=\scriptsize, color=inkprimary, anchor=north west] at (axis cs:2.03,46) {48.4, $\times$1.70};
\node[font=\scriptsize, color=inkprimary, anchor=north] at (axis cs:3,56.8) {59.8, $\times$2.10};
\end{groupplot}
\end{tikzpicture}
\caption{Books indexed per second over a whole rebuild of the 2{,}800 books, mean and 95\,\% confidence interval of 5
repetitions. Left: batch size with one replica. Right: replicas with batches of 100, against linear scaling from one
replica. \run{20261105-220915-index-throughput-c1376c7}.}
\label{fig:index}
\end{figure}

\begin{itemize}
\item \textbf{Bigger batches pay off, then flatten} (Figure~\ref{fig:index}), as expected: $\times$2.3 from 1 to 10
  books, $\times$2.3 from 10 to 100, and only +17\,\% from 100 to 250. Each batch pays a fixed cost (its lease,
  publication and marking) that more books share.
\item \textbf{Replicas help less than linearly}, as expected: $\times$1.70 with 2 and $\times$2.10 with 3, that is 85\,\%
  and 70\,\% of linear.
\item \textbf{Batches of 250 fit in a 512\,MB heap}, as expected: no indexer ran out of memory, and the container
  peaked at 745--808\,MiB per trial with batches of 250, against 587--778\,MiB with batches of 100.
\item \textbf{MongoDB is where the time goes.} JMH writes a batch of 100 books in 2.81\,s, which alone caps one
  indexer at 36 books/s, close to the 28.5 measured end to end. Counting the terms of the largest book takes 37\,ms.
  During indexing MongoDB's CPU grows faster than the indexers' (Table~\ref{tab:resources}).
\end{itemize}

\textbf{Indexing latency.} Over a whole rebuild, each book costs the inverse of the throughput: 0.18\,s with batches of
1, 0.035\,s with 100 and 0.017\,s with three replicas of 100. A single book also waits for its batch: from its claim
to its batch's publication it takes 3.45\,s with one replica and 4.93\,s with three (Table~\ref{tab:coordination}).

\subsection{Search latency under concurrency}
\label{sec:search}

\begin{figure}[ht]
\centering
\begin{tikzpicture}
\begin{groupplot}[
  group style={group size=2 by 1, horizontal sep=1.7cm},
  report,
  width=7.6cm, height=5.2cm,
  xmode=log, log basis x=2,
  xtick={1,8,32,64}, xticklabels={1,8,32,64},
  xlabel={Closed-loop clients},
  ymin=0,
  clip=false,
  error bars/y dir=both, error bars/y explicit,
]
\nextgroupplot[ylabel={Requests per second}, ymax=5000, title={Throughput}, yticklabel style={/pgf/number format/1000 sep={{,}}}]
\addplot[color=modehazelcast, mark=*, mark size=2.4pt, error bars/error bar style={line width=0.7pt}]
  table[x=clients, y=rps, y error=rps_ci, restrict expr to domain={\thisrow{threads}*10+\thisrow{replicas}}{321:321}] {\searchresults};
\node[font=\scriptsize, color=inkprimary, anchor=south east] at (axis cs:1.15,700) {646};
\node[font=\scriptsize, color=inkprimary, anchor=north west] at (axis cs:8.4,3550) {3{,}618};
\node[font=\scriptsize, color=inkprimary, anchor=south east] at (axis cs:64,4520) {4{,}374 and 4{,}381};
\nextgroupplot[ylabel={Latency (ms)}, ymax=30, xmax=110, title={Latency}, legend style={at={(0.03,0.97)}, anchor=north west}]
\addplot[color=modehazelcast, mark=*, mark size=2.4pt, error bars/error bar style={line width=0.7pt}]
  table[x=clients, y=mean_ms, y error=ci_ms, restrict expr to domain={\thisrow{threads}*10+\thisrow{replicas}}{321:321}] {\searchresults};
\addlegendentry{mean}
\addplot[color=modeactivemq, mark=square*, mark size=2.2pt, dashed]
  table[x=clients, y=p95_ms, restrict expr to domain={\thisrow{threads}*10+\thisrow{replicas}}{321:321}] {\searchresults};
\addlegendentry{p95}
\node[font=\scriptsize, color=inkprimary, anchor=east] at (axis cs:60,27.2) {27.2};
\node[font=\scriptsize, color=inkprimary, anchor=west] at (axis cs:66,14.5) {14.5};
\end{groupplot}
\end{tikzpicture}
\caption{\code{GET /search} with one \code{query} replica and 32 HTTP threads, as the closed-loop clients grow. Mean and
95\,\% confidence interval of 5 repetitions of 120\,s each, after 30\,s of warmup; the p95 is the mean of each
repetition's p95. \run{20261104-140900-search-latency-8e1c1e8}.}
\label{fig:search}
\end{figure}

\begin{table}[ht]
\centering
\small
\caption{\code{GET /search} by clients, HTTP threads and \code{query} replicas, cycling through
\code{workload/queries.txt}; 18{,}496{,}487 requests in all. \run{20261104-140900-search-latency-8e1c1e8}.}
\label{tab:search-latency}
\begin{tabular}{@{}rrrrrr@{}}
\toprule
Clients & HTTP threads & Replicas & Mean latency (ms) & p95 (ms) & Requests/s \\
\midrule
1 & 32 & 1 & 1.54 $\pm$ 0.02 & 2.73 & 646 $\pm$ 8 \\
8 & 32 & 1 & 2.20 $\pm$ 0.06 & 3.87 & 3{,}618 $\pm$ 99 \\
32 & 32 & 1 & 7.29 $\pm$ 0.04 & 14.5 & 4{,}374 $\pm$ 26 \\
64 & 32 & 1 & 14.5 $\pm$ 0.2 & 27.2 & 4{,}381 $\pm$ 50 \\
32 & 8 & 1 & 12.9 $\pm$ 1.3 & 19.3 & 2{,}493 $\pm$ 254 \\
32 & 128 & 1 & 7.56 $\pm$ 0.06 & 15.1 & 4{,}217 $\pm$ 36 \\
32 & 32 & 2 & 5.66 $\pm$ 0.02 & 10.1 & 5{,}625 $\pm$ 24 \\
32 & 32 & 3 & 5.82 $\pm$ 0.04 & 10.5 & 5{,}471 $\pm$ 39 \\
\bottomrule
\end{tabular}
\end{table}

\begin{itemize}
\item \textbf{One replica saturates at about 4{,}400 requests/s} with 32 clients (Figure~\ref{fig:search}). From 32
  to 64 clients the throughput stays the same and the latency doubles, which is Little's law~\cite{little} for a saturated server:
  64 clients / 4{,}381 requests/s = 14.6\,ms, the measured 14.5\,ms.
\item \textbf{Threads must cover the clients.} With 8 HTTP threads and 32 clients, requests queue and latency is
  77\,\% higher than with 32 threads; 128 threads add nothing and lose 4\,\% of the requests/s
  (Table~\ref{tab:search-latency}).
\item \textbf{A second replica helps, a third does not.} Two replicas give +29\,\% requests/s and $-22$\,\% latency;
  three give the same as two. With two replicas the limit is no longer the \code{query} JVM but what the replicas
  share: nginx, MongoDB and the host's CPUs, which the load generator also uses.
\item \textbf{The ranking is not the cost.} JMH ranks a three-word query in 60\,\textmu s, about 1\,\% of the 5--7\,ms
  each request takes. The rest is reading postings and metadata from MongoDB, HTTP and JSON.
\item \textbf{Answers stay correct.} Every request answered \code{200} except 170 \code{503}s, 168 of them with 64
  clients: 0.006\,\% of that configuration's 2{,}628{,}660 requests.
\end{itemize}

\subsection{Coordination cost}
\label{sec:coordcost}

\begin{table}[ht]
\centering
\small
\caption{Where a book's time goes in mode \mode{hazelcast}, by replicas of \code{crawler} and \code{indexer} together:
the wait in \code{books-to-download} and in \code{books-to-index}, the work, and the time from enqueue to publication.
Mean of 5 repetitions $\pm$ 95\,\% interval,
from the transition logs. \mode{hazelcast} trials of \run{20261106-012054-coordination-variants-c1376c7}.}
\label{tab:coordination}
\footnotesize
\setlength{\tabcolsep}{4pt}
\begin{tabular}{@{}rrrrrr@{}}
\toprule
N & Wait for a crawler & Download & Wait for an indexer & Indexing & Enqueue to publication \\
\midrule
1 & 1{,}383 $\pm$ 327 ms & 12.6 $\pm$ 1.6 ms & 46.6 $\pm$ 1.9 s & 3.45 $\pm$ 0.13 s & 51.4 $\pm$ 1.9 s \\
2 & 670 $\pm$ 115 ms & 27.0 $\pm$ 3.8 ms & 27.0 $\pm$ 2.7 s & 4.19 $\pm$ 0.41 s & 31.9 $\pm$ 3.1 s \\
3 & 736 $\pm$ 127 ms & 44.4 $\pm$ 4.5 ms & 19.9 $\pm$ 0.7 s & 4.93 $\pm$ 0.21 s & 25.6 $\pm$ 0.9 s \\
\bottomrule
\end{tabular}
\end{table}

\begin{itemize}
\item \textbf{Coordination costs far more time than the work}, as expected (Table~\ref{tab:coordination}): with one
  replica a book waits 1.4\,s for a crawler to download it in 13\,ms.
\item \textbf{More consumers on one queue help little.} The wait for a crawler halves with 2 replicas and does not fall
  further with 3, while the download itself gets slower with every replica, 13, 27 and 44\,ms.
\item \textbf{Indexing is the bottleneck of the pipeline.} The wait in \code{books-to-index} is 91\,\% of the time to
  publication with one replica. It falls 42\,\% with 2 replicas and 57\,\% with 3; from 1 to 3 replicas that is more
  than the 47\,\% of the download wait, against the hypothesis. Books wait for an indexer, not for their batch to fill: with one replica the p95 of that wait is
  87\,s, far above the 30\,s a batch waits at most.
\end{itemize}

\subsection{Resource usage}
\label{sec:resources}

\begin{table}[ht]
\centering
\scriptsize
\setlength{\tabcolsep}{4pt}
\caption{CPU, in percent of one core as \code{docker stats} counts it, and memory during each phase, by replicas of
\code{crawler} and \code{indexer}. ``Stack'' adds up every container; ``Host'' is the whole machine.
\run{20261105-235758-resource-usage-c1376c7}.}
\label{tab:resources}
\begin{tabular}{@{}rlrrrrrr@{}}
\toprule
N & Phase & Crawlers' CPU & Indexers' CPU & MongoDB's CPU & Stack CPU & Stack memory (MiB) & Host CPU \\
\midrule
1 & Ingestion & 378 $\pm$ 94\,\% & -- & 50 $\pm$ 9\,\% & 505 $\pm$ 105\,\% & 1{,}199 $\pm$ 122 & 956 $\pm$ 85\,\% \\
2 & Ingestion & 512 $\pm$ 74\,\% & -- & 62 $\pm$ 11\,\% & 680 $\pm$ 109\,\% & 1{,}398 $\pm$ 125 & 1{,}163 $\pm$ 155\,\% \\
3 & Ingestion & 537 $\pm$ 53\,\% & -- & 65 $\pm$ 4\,\% & 717 $\pm$ 32\,\% & 1{,}723 $\pm$ 73 & 1{,}281 $\pm$ 71\,\% \\
1 & Indexing & -- & 64 $\pm$ 2\,\% & 77.9 $\pm$ 0.4\,\% & 147 $\pm$ 2\,\% & 2{,}264 $\pm$ 199 & 208 $\pm$ 7\,\% \\
2 & Indexing & -- & 119 $\pm$ 10\,\% & 150 $\pm$ 3\,\% & 277 $\pm$ 7\,\% & 2{,}967 $\pm$ 148 & 346 $\pm$ 12\,\% \\
3 & Indexing & -- & 176 $\pm$ 11\,\% & 220 $\pm$ 3\,\% & 406 $\pm$ 10\,\% & 3{,}930 $\pm$ 134 & 494 $\pm$ 11\,\% \\
\bottomrule
\end{tabular}
\end{table}

\begin{itemize}
\item \textbf{Crawlers share a fixed rate.} Their total CPU grows with replicas, but each crawler uses less, 378, 256
  and 179\,\%, instead of staying flat as expected: the same books per second are split among more crawlers, and the
  24 hardware threads are never near saturation.
\item \textbf{MongoDB follows the indexers.} During ingestion the crawlers use 7.6--8.3 times MongoDB's CPU; during
  indexing MongoDB uses more CPU than the indexers and grows with them, from 78\,\% to 220\,\%, as expected.
\item \textbf{Memory is bounded.} A crawler or indexer container reaches 497--683\,MiB at p99, above its
  \code{-Xmx512m} in all but one configuration, because \code{docker stats} counts the whole JVM (heap, metaspace,
  thread stacks and buffers), not the heap. No container grows with the load; \code{control} and \code{query} stay
  below 270\,MiB.
\item \textbf{The host pays for Docker.} The host uses 1.7--1.9 times the stack's CPU during ingestion and 1.2--1.4
  times during indexing: the rest runs in Docker Desktop's VM and Windows.
\end{itemize}

\subsection{Coordination mechanisms compared}
\label{sec:modes}

\code{coordination-variants} ran the whole corpus through the five modes of Table~\ref{tab:modes}, with 1, 2 and 3
replicas of \code{crawler} and \code{indexer} together, 5 interleaved repetitions each: 75 trials. Its
\code{pairs.csv} holds the difference of every pair of modes, for every metric and replica count, with its 95\,\%
confidence interval. Figure~\ref{fig:modes} shows the pipeline metrics, and Table~\ref{tab:decision} in
Section~\ref{sec:winner} the rest.

\begin{figure}[ht]
\centering
\begin{tikzpicture}
\begin{groupplot}[
  group style={group size=2 by 1, horizontal sep=1.7cm},
  report,
  width=7.6cm, height=5.6cm,
  xtick={1,2,3}, xmin=0.75, xmax=3.25,
  xlabel={Replicas of \code{crawler} and \code{indexer}},
  ymin=0,
]
\nextgroupplot[ylabel={Books per second}, ymax=70, title={Pipeline throughput}, legend to name=modeslegend, legend columns=5, legend style={/tikz/every even column/.append style={column sep=0.35cm}}]
\addplot[color=modehazelcast, mark=*, modeline] table[x expr=\thisrow{replicas}-0.08, y=hz, y error=hzci] {\modesthroughput};
\addlegendentry{\mode{hazelcast}}
\addplot[color=modeactivemq, mark=square*, modeline] table[x expr=\thisrow{replicas}-0.04, y=amq, y error=amqci] {\modesthroughput};
\addlegendentry{\mode{activemq}}
\addplot[color=modepersistent, mark=diamond*, modeline, mark size=3pt] table[x expr=\thisrow{replicas}, y=amqp, y error=amqpci] {\modesthroughput};
\addlegendentry{\mode{activemq-persistent}}
\addplot[color=modepolling, mark=triangle*, modeline, mark size=3pt] table[x expr=\thisrow{replicas}+0.04, y=poll, y error=pollci] {\modesthroughput};
\addlegendentry{\mode{polling}}
\addplot[color=modeorchestrator, mark=pentagon*, modeline, mark size=2.8pt] table[x expr=\thisrow{replicas}+0.08, y=orch, y error=orchci] {\modesthroughput};
\addlegendentry{\mode{orchestrator}}
\node[font=\scriptsize, color=inkprimary, anchor=north east] at (axis cs:3.2,33) {\mode{orchestrator}: 34.6};
\nextgroupplot[ylabel={Mean time to publication (s)}, ymax=62, title={Time from enqueue to publication}]
\addplot[color=modehazelcast, mark=*, modeline] table[x expr=\thisrow{replicas}-0.08, y=hz, y error=hzci] {\modespublication};
\addplot[color=modeactivemq, mark=square*, modeline] table[x expr=\thisrow{replicas}-0.04, y=amq, y error=amqci] {\modespublication};
\addplot[color=modepersistent, mark=diamond*, modeline, mark size=3pt] table[x expr=\thisrow{replicas}, y=amqp, y error=amqpci] {\modespublication};
\addplot[color=modepolling, mark=triangle*, modeline, mark size=3pt] table[x expr=\thisrow{replicas}+0.04, y=poll, y error=pollci] {\modespublication};
\addplot[color=modeorchestrator, mark=pentagon*, modeline, mark size=2.8pt] table[x expr=\thisrow{replicas}+0.08, y=orch, y error=orchci] {\modespublication};
\end{groupplot}
\node at ($(group c1r1.south)!0.5!(group c2r1.south)+(0,-1.25cm)$) {\pgfplotslegendfromname{modeslegend}};
\end{tikzpicture}
\caption{The same pipeline in each coordination mode: books published per second over the whole corpus, and mean time
from a book's enqueue to the publication of its batch. Mean and 95\,\% confidence interval of 5 repetitions; the modes
are shifted slightly along the x axis so their intervals do not overlap.
\run{20261106-012054-coordination-variants-c1376c7}.}
\label{fig:modes}
\end{figure}

The script stated one hypothesis per mode before measuring. The common hypothesis, that the work per book is the same
in every mode, holds for indexing but not for downloads. In \mode{activemq-persistent} and \mode{orchestrator} a
download takes 8--15\,ms at p50 with 2 and 3 replicas, against 21--45\,ms elsewhere, because books reach the crawlers
more slowly than the crawlers take them, so fewer downloads overlap. That changes milliseconds against a time to
publication of tens of seconds.

\begin{itemize}
\item \textbf{\mode{activemq}} has the throughput and time to publication of \mode{hazelcast} at every N, within
  their intervals, as expected, with 1.5 times its coordinator CPU.
\item \textbf{\mode{activemq-persistent}} has the throughput of \mode{activemq} at every N, as expected, with 1.3--1.4
  times its broker CPU.
\item \textbf{\mode{polling}} has the throughput of \mode{hazelcast} at every N and 10--12\,\% more MongoDB CPU, from
  the sorted claims and the idle polls. Visibility at p50 is 2.3--2.7\,s, about half the 5\,s refresh period. All three
  matched the hypotheses.
\item \textbf{\mode{orchestrator}} keeps the throughput with 1 and 2 replicas and loses 41\,\% with 3, which its
  hypothesis allowed, because one thread sends one \code{POST} at a time. Its mean time to publication is 4.3--6.3\,s
  above \mode{hazelcast} (Figure~\ref{fig:modes}), and its median 4.7--5.5\,s (Table~\ref{tab:decision}), not the two
  ticks expected: books wait in MongoDB until \code{control} dispatches them, 4.2--5.6\,s on average against
  0.7--1.4\,s in \code{books-to-download}. Its visibility is 168--211\,ms at p50, under the expected 300\,ms.
\item \textbf{Coordination is cheap in every mode.} The broker or \code{control} uses at most 17\,\% of one core,
  against MongoDB's 80--247\,\%.
\end{itemize}

\textbf{Publication visibility.} The time from a batch's publication to the moment a \code{query} replica holds it
in its snapshot depends on how \code{query} learns of the batch (Table~\ref{tab:visibility}). With the topic listener
a batch is searchable in about 40\,ms, under the expected 100\,ms. Without it, delays spread between 0 and the 5\,s
period, with a median of 2.7\,s. The listener costs no CPU that the run can see: each \code{query} replica uses about
3\,\% of a core either way.

\begin{table}[ht]
\centering
\small
\caption{Publication visibility with 2 \code{query} replicas, one value per batch and replica.
\run{20261106-002208-publication-visibility-c1376c7}.}
\label{tab:visibility}
\begin{tabular}{@{}lrrrr@{}}
\toprule
\code{published-batches} listener & Mean (ms) & p50 (ms) & p95 (ms) & p99 (ms) \\
\midrule
On & 43.6 $\pm$ 7.3 & 37 & 54 & 145 \\
Off & 2{,}605 $\pm$ 251 & 2{,}672 & 4{,}880 & 4{,}965 \\
\bottomrule
\end{tabular}
\end{table}

\textbf{Broker latency in isolation.} \code{broker-latency} sends 30{,}000 messages per trial through each broker, with
1, 3 and 8 consumers (Figure~\ref{fig:broker}). Every p50 is between 0.1 and 0.2\,ms: 102--105\,\textmu s for Hazelcast,
137--139\,\textmu s for ActiveMQ and 180--185\,\textmu s for persistent ActiveMQ. Persistence costs 1.3 times the latency,
not the 3 times expected, but it caps the throughput at about 1{,}460 messages/s whatever the consumers, because each
send waits for its disk write~\cite{activemq}. Non-persistent ActiveMQ scales with consumers, $\times$2.8 from 1 to 8, while Hazelcast
stays at about 8{,}700 messages/s. A message is never what a book waits for: the slowest p50, 185\,\textmu s, is about
7{,}700 times shorter than the median wait in \code{books-to-download} with one replica.

\begin{figure}[ht]
\centering
\begin{tikzpicture}
\begin{axis}[
  report,
  width=12cm, height=5.2cm,
  ybar, bar width=0.42cm,
  symbolic x coords={1,3,8}, xtick=data,
  enlarge x limits=0.22,
  xlabel={Consumers},
  ylabel={Messages per second},
  ymin=0, ymax=32000,
  scaled y ticks=false,
  yticklabel style={/pgf/number format/fixed, /pgf/number format/1000 sep={{,}}},
  legend style={at={(0.02,0.97)}, anchor=north west},
  error bars/y dir=both, error bars/y explicit,
  every axis plot/.append style={draw=none, error bars/error bar style={draw=inkprimary, line width=0.6pt}},
]
\addplot[fill=modehazelcast] table[x=consumers, y=hz, y error=hzci] {\brokerresults};
\addlegendentry{Hazelcast}
\addplot[fill=modeactivemq] table[x=consumers, y=amq, y error=amqci] {\brokerresults};
\addlegendentry{ActiveMQ, non-persistent}
\addplot[fill=modepersistent] table[x=consumers, y=amqp, y error=amqpci] {\brokerresults};
\addlegendentry{ActiveMQ, persistent}
\end{axis}
\end{tikzpicture}
\caption{Throughput of each broker in isolation, mean and 95\,\% confidence interval of 5 repetitions of 30{,}000
messages: Hazelcast 8{,}513--8{,}872 messages/s, non-persistent ActiveMQ 9{,}350, 18{,}338 and 25{,}904, persistent
ActiveMQ 1{,}454--1{,}475. \run{20261106-123534-broker-latency-b4667d8}.}
\label{fig:broker}
\end{figure}

\subsection{Recovery}

Each of the 165 recovery runs causes one failure of the matrix (Appendix~\ref{app:failures}) and measures the time
until every book is \code{PUBLISHED}, found exactly once by \code{/search} and in a single published batch
(Table~\ref{tab:recovery}). The runs use the short leases of \code{scripts/recovery/recovery.env} (6\,s) and run once
per row, mode and replica count, so they have no confidence interval. \code{recovery-compare} checks that the runs of a
row give the same books and search answers whatever the replicas and mode, and it passed for all 14 rows.

\begin{table}[ht]
\centering
\small
\caption{Seconds from the failure to correct answers. Rows 09--12 seed the whole mirror, so their time includes
processing the 2{,}800 books; the others seed the 20 sample books. Runs \code{2026110[56]-...-recovery-<nn>[-<mode>]-r<N>-c1376c7},
summarized by \run{20261106-073908-recovery-summary-c1376c7}.}
\label{tab:recovery}
\begin{tabularx}{\textwidth}{@{}lYrL{3.5cm}L{3.9cm}@{}}
\toprule
Rows & Failure & Corpus & \mode{hazelcast}, N = 1/2/3 & Other modes \\
\midrule
01--04 & \code{crawler} crash points & sample & 8.6--18.8 & 7.2--22.8 \\
05--08 & \code{indexer} crash points & sample & 6.3--17.7 & 6.6--16.9 \\
09 & \code{crawler} paused longer than its lease & mirror & 373 / 276 / 225 & 363--373, N = 1 \\
10 & \code{indexer} paused longer than its lease & mirror & 382 / 278 / 241 & 368--441, N = 1 \\
11 & \code{control} killed while seeding & mirror & 367 / 255 / 253 & 365--373, N = 1 \\
12 & Broker restarted & mirror & 374 / 261 / 235 & 372 and 384, ActiveMQ modes, N = 1 \\
13 & MongoDB stopped for 10\,s & sample & 16.6 / 15.9 / 14.8 & 14.7--22.6 \\
14 & A \code{query} replica killed & sample & -- / 0.9 / 0.7 & 0.2--1.2 \\
\bottomrule
\end{tabularx}
\end{table}

Every row recovers in every mode where it applies, with the same answers. The crash points of \code{crawler} and
\code{indexer} recover within 23\,s with the 20 sample books. Rows 09--12 take 225--441\,s, mostly ingesting and
indexing the corpus, and that time falls with N as indexing gets faster. Killing one of two or three \code{query}
replicas costs at most 1.2\,s of search answers, because nginx sends the requests to the others.

\textbf{How lost work comes back.} \code{recovery-delay} reads the 29 runs of rows 04, 05 and 12, where a crash or a
broker restart loses work after its message was taken, and measures for each book the time from the failure to its
requeue and to its next claim (Table~\ref{tab:delay}). In \mode{hazelcast} and \mode{activemq}, \code{control} requeues
the books of rows 04 and 05 11.7--11.8\,s after the failure at the median, and every one of them within
10.8--11.8\,s, as expected. \mode{activemq-persistent} requeues nothing: the broker redelivers, and books are claimed
again in 3.8--6.0\,s at the median. The orchestrator dispatches again 9.7\,s after the failure at the median, and every
book within 9.5--10.1\,s, because its 10\,s count from the book's last dispatch, which comes before the crash.

\textbf{Row 12 loses two kinds of messages.} Download messages come back like those of rows 04 and 05: the median
requeue of each run is 11.7--12.7\,s after the restart, whatever N. Index messages wait a median of 207, 97 and 82\,s
in the runs with 1, 2 and 3 replicas, because \code{control} only requeues into an empty queue (\spec{5.3}, step~4).
After the restart the crawlers keep enqueuing the books they download, so \code{books-to-index} only empties once the
indexers, with the batches of 5 of the recovery tests, have taken every one of them. The transition lines of the
three runs show it: the requeue comes in the first reconciliation after the indexers take the last message of the
queue. The wait therefore follows indexing, and the more indexers the shorter it is. In \mode{activemq-persistent}
neither kind is lost, and books are claimed again 6.0 and 11.7\,s after the restart at the median.

\begin{table}[ht]
\centering
\small
\caption{Median seconds from the failure to a book's requeue (or new dispatch) and to its next claim, over the books
of every run of the row, with the mean of the two middle values when their number is even; row 12 apart for the books
that were waiting for a download and for an index. \run{20261106-073911-recovery-delay-c1376c7}.}
\label{tab:delay}
\setlength{\tabcolsep}{4.5pt}
\begin{tabular}{@{}lrrrr@{}}
\toprule
& \multicolumn{2}{c}{Rows 04 / 05} & \multicolumn{2}{c}{Row 12, download / index messages} \\
\cmidrule(lr){2-3}\cmidrule(l){4-5}
Mode & Requeue & Claim & Requeue & Claim \\
\midrule
\mode{hazelcast} & 11.8 / 11.7 & 11.5 / 5.6 & 11.9 / 96.8 & 9.3 / 134.3 \\
\mode{activemq} & 11.8 / 11.7 & 10.8 / 8.5 & 11.7 / -- & 13.5 / -- \\
\mode{activemq-persistent} & none & 3.8 / 5.0 & none & 6.0 / 11.7 \\
\mode{orchestrator} & 9.7 / 9.7 & 5.0 / 9.8 & -- & -- \\
\bottomrule
\end{tabular}
\end{table}

\subsection{Bottlenecks and the limits of this evidence}
\label{sec:limits}

\begin{enumerate}
\item \textbf{Indexing limits the pipeline.} One indexer publishes 28.5 books/s, while one crawler downloads about 500.
  Books therefore wait 20--47\,s in \code{books-to-index}, and an indexer spends its time writing postings to the one
  MongoDB.
\item \textbf{Ingestion stops at 500--670 books/s} with 8 or more workers, whatever the replicas, with at most 53\,\% of
  the host's CPU in use. The download workers wait on the file I/O they share, the mirror and the datalake volume; the
  queue and MongoDB take under 3\,ms per book.
\item \textbf{Search saturates one \code{query} replica at about 4{,}400 requests/s.} A second replica adds 29\,\%,
  and a third adds nothing.
\item \textbf{Coordination is not a bottleneck in any mode.} A message takes 0.1--0.2\,ms, and the broker or
  \code{control} uses less than 17\,\% of a core. The mode changes how soon a book is searchable, at the median:
  38--48\,ms with a listener, 168--211\,ms with the orchestrator, and 2.3--2.7\,s with polling
  (Table~\ref{tab:decision}).
\end{enumerate}

These results describe one host. Every replica shares its CPUs, its disk, and the one MongoDB and broker, so a speedup
here shows how the design uses one machine, not how it will scale across several. The evidence has other limits too:
\begin{itemize}
\item \code{search-latency} and the indexer's JMH come from 3 and 4 November, reused because their code did not change.
\item After the campaign, \code{dev} changed one class of \code{control} (\code{65095a7}): the orchestrator now asks
  \code{query} to refresh once per tick with new batches, not once per batch. Only the visibility of mode
  \mode{orchestrator} can differ from what was measured at \code{c1376c7}.
\item A code review after the campaign fixed eleven defects (\code{SPEC.md} 1.11). One touches a measured path: in both
  \mode{activemq} modes, every queue size \code{control} read, two per reconciliation and two per
  \code{GET /control/status}, which \code{bench wait} asks every second, left a temporary queue in the broker until the
  trial ended. Their coordinator CPU includes that extra work, which was not measured on its own, so it is an upper
  bound.
\item The campaign left an idle ActiveMQ container running during the recovery runs of \mode{polling} and
  \mode{orchestrator}. No service of those modes connects to it, and the scripts bring every mode down since
  \code{b4667d8}.
\item Docker Desktop on Windows adds a VM: the host uses 1.2--1.9 times the CPU of the containers. \code{docker stats}
  reports a container's whole memory, not the JVM heap.
\item An ingestion phase lasts about 5\,s, about 10 \code{docker stats} samples per repetition, so its CPU means are
  coarse. Recovery runs have one sample per row, mode and replica count, and in the four alternative modes rows 09--12
  only ran with one replica.
\end{itemize}
\section{Conclusions and Future Improvements}
\label{sec:conclusions}

\subsection{What the measurements say}

\textbf{Strengths.} The design is correct under every failure we could cause. All 165 recovery runs recover with the
same answers whatever the mode and the replica count, and no book ever appeared twice. Coordination costs little: a
message takes 0.1--0.2\,ms, and the broker uses at most 17\,\% of one core. Indexer and \code{query} replicas raise
throughput on one host, $\times$2.10 with three indexers and +29\,\% with a second \code{query}. A book becomes
searchable about 40\,ms after its batch is published. The contracts held: the API checks pass in every mode and
with one or three \code{query} replicas, and every figure traces to a run.

\textbf{Weaknesses.} Indexing is the bottleneck: one indexer publishes 28.5 books/s against about 500 downloaded, and
its time goes into MongoDB writes. Crawler replicas add nothing beyond 500--670 books/s, because their workers wait on
the file I/O they all share, not on the queue or MongoDB. \code{control} is a single process that every mode needs,
since it alone releases leases and drives rebuilds. When the broker loses messages, \code{control} brings them back
after \code{TARANTINO_REQUEUE_SECONDS}, slower than the redelivery of a persistent broker, and only into an empty queue:
after a Hazelcast restart, lost index messages waited 82--207\,s for \code{books-to-index} to empty.

\subsection{From one host to the Stage~3 cluster}
\label{sec:stage3}

The conditional writes already work across machines. Claims, leases and expirations compare with MongoDB's clock,
never a host's, and no process holds a lock that another needs. What Stage~3 has to replace is everything named after
the host (Table~\ref{tab:stage3}): Compose service names and nginx's resolver \code{127.0.0.11}, the local
\code{datalake} volume, the single Hazelcast member and MongoDB server, and, in mode \mode{orchestrator}, the JVM's DNS
cache through which \code{control} finds the \code{query} replicas.

\begin{table}[ht]
\centering
\small
\caption{Each measured limit, the single-host assumption behind it, and the change Stage~3 makes.}
\label{tab:stage3}
\begin{tabularx}{\textwidth}{@{}L{4.0cm}L{4.1cm}Y@{}}
\toprule
Measured limit & Single-host cause & Stage~3 change \\
\midrule
Indexing: 28.5 books/s per indexer, $\times$2.10 with 3, with MongoDB's CPU growing faster than the indexers' & Every indexer writes postings to one MongoDB server & Postings sharded by \code{term} across a sharded MongoDB, so indexers on several nodes write to different shards; the batch, its lease and its publication stay single-document writes \\
Ingestion: 500--670 books/s whatever the replicas, with the CPU far from full & Every crawler reads the one mirror and writes the one datalake volume of the host & Crawlers on several nodes, each with its own disk: the datalake split by \code{book_id} across the nodes, with each book's node in \code{book_states}, so file I/O grows with the nodes \\
Search: one \code{query} replica saturates at about 4{,}400 requests/s, and a third adds nothing to two & Replicas share nginx, the one MongoDB and the host's CPUs with the load generator & \code{query} replicas on several nodes behind a load balancer, reading a MongoDB replica set; the snapshot already lets each replica answer on its own \\
Recovery: lost messages wait for \code{control}'s requeue & \code{control} is one process, acceptable on one host & \code{control} on several nodes with one leader, elected through a lease document in MongoDB; lost work acknowledged after the job (below) \\
Mode \mode{orchestrator}: 41\,\% less throughput with 3 replicas & One thread dispatches every book, one \code{POST} at a time & Not carried to Stage~3 \\
\bottomrule
\end{tabularx}
\end{table}

\textbf{Communication and storage across machines.} Services keep talking only through the broker and MongoDB, so
moving them to other nodes changes addresses, not code. Hazelcast becomes a cluster of three members, where each queue
item has a synchronous backup on another member, so one member can fail without losing queued work. MongoDB becomes a
replica set, with majority writes for every conditional write, so a failover never loses a claim. The datalake moves to
storage that every node mounts.

\textbf{Distribution of tasks and requests.} Crawlers and indexers keep competing for queues, which needs no
coordination between them. Partitioning \code{books-to-download} and \code{books-to-index} by \code{book_id} spreads
them over the members, so no member holds every message. Search requests are spread by a load balancer, as nginx does
now.

\textbf{Consistency and unavailable nodes.} Publication stays one conditional write, and \code{query} keeps reading a
snapshot of published batches, so a reader never sees a half-built index on any node. A node that disappears leaves
leases that expire, and \code{control} releases them exactly as it does after a crash today: rows 01 to 10 of the
failure matrix already exercise that path. A network partition cannot create two owners of a book, because ownership
is a MongoDB document, not a broker message.

\subsection{The winning mode in Stage~3}

\mode{hazelcast} wins under the decision rule fixed before the campaign (Section~\ref{sec:winner}), so Stage~3 keeps it
and does not migrate to another broker. The runner-up, \mode{activemq-persistent}, showed the one property worth
taking: lost work comes back without \code{control}, in 3.8--6.0\,s instead of about 12\,s, because the message is
acknowledged only after the job. Stage~3 adds that property to the Hazelcast design and measures it with the same
rows and the same \code{recovery-delay} experiment:
\begin{enumerate}
\item Take each book inside a Hazelcast transaction that commits after the claim and the job, so a consumer that dies
  returns its message to the queue. Hazelcast's transactional queue offers that rollback~\cite{hazelcast}, and the claim in MongoDB
  still drops any duplicate it produces.
\item Requeue lost messages without waiting for their queue to empty, for example as soon as a queue holds fewer
  messages than the books that wait for it, and measure row 12 again.
\item Elect the \code{control} leader through a lease in MongoDB, so reconciliation survives the loss of a node.
\item Repeat \code{coordination-variants} for \mode{hazelcast} and \mode{activemq-persistent} on the cluster. If the
  transactional queue does not close the gap in rows 04, 05 and 12, persistent ActiveMQ replaces Hazelcast as the
  broker. The rest of the code stays the same, because both already sit behind the same ports.
\end{enumerate}

A smaller improvement follows from this report: batches of 250 books once the indexer has a larger heap, which gave
17\,\% more throughput. The rebuild race of Section~\ref{sec:indexer} is already closed in \code{SPEC.md} 1.11.

\subsection{Lessons}

\begin{itemize}
\item \textbf{Measure the alternatives.} Arguing against orchestration and polling was easy; measuring them showed
  where each one actually loses, by how much, and what the runner-up does better.
\item \textbf{Keep the truth in one place.} Treating every message as a hint and every transition as one conditional
  write made 14 different failures recoverable by the same reconciliation.
\item \textbf{Keep the evidence with the results.} Five interleaved repetitions, hypotheses written before measuring,
  and a versioned run behind every figure made every number in this report checkable against its raw data, and we
  checked them all before submitting it.
\end{itemize}
\section{Statement on AI Usage}
\label{sec:ai}

In accordance with the course guidelines, we disclose the use of AI tools during this stage, mainly Claude
(Anthropic) through Claude Code. AI was used as an assistant to speed up routine work. The architecture, the contracts
in \code{SPEC.md}, the design of every experiment, the decision rule for the coordination modes and the interpretation
of the results were decided and reviewed by the team. Specific uses included:

\begin{itemize}
\item \textbf{Code and tests.} Drafting and reviewing implementation code, tests and scripts against \code{SPEC.md},
  including parts of the alternative coordination modes. Every change went through a pull request and the five
  required checks of continuous integration before it reached \code{dev}.
\item \textbf{Reviews and audits.} Checking the repository against the assignment and the specification, and finding
  inconsistencies between the documents, the code and the results.
\item \textbf{Data processing and writing.} Extracting tables from the run files, checking every figure in this report
  against the raw data, and drafting and formatting the documentation and this report in \LaTeX.
\end{itemize}

Every piece of AI-assisted code or text was reviewed, tested and integrated by the team member who owns it, to keep
it consistent with the specification and with academic integrity.
\begin{thebibliography}{99}
\addcontentsline{toc}{section}{References}

\bibitem{assignment}
Big Data, Universidad de Las Palmas de Gran Canaria.
\newblock \emph{Stage 2: Building a Service-Oriented Architecture}. Search Engine Project, 2026--2027.

\bibitem{stage1}
Query Tarantino.
\newblock \emph{Search Engine Project: Stage 1, Building the Data Layer}. Report, 2026.

\bibitem{rfc9457}
M.~Nottingham, E.~Wilde and S.~Dalal.
\newblock \emph{Problem Details for HTTP APIs}. RFC 9457, IETF, 2023.

\bibitem{openapi}
OpenAPI Initiative.
\newblock \emph{OpenAPI Specification}, version 3.1.0, 2021.

\bibitem{mongodb}
MongoDB, Inc.
\newblock \emph{MongoDB 7.0 Manual}: updates with aggregation pipelines, the \code{$$NOW} variable and \code{findOneAndUpdate}.
\newblock \url{https://www.mongodb.com/docs/v7.0/}.

\bibitem{hazelcast}
Hazelcast, Inc.
\newblock \emph{Hazelcast Platform 5.7 Documentation}: queues, topics, backups and transactions.
\newblock \url{https://docs.hazelcast.com/hazelcast/5.7/}.

\bibitem{activemq}
The Apache Software Foundation.
\newblock \emph{Apache ActiveMQ Classic 6 Documentation}: persistence and message acknowledgement.
\newblock \url{https://activemq.apache.org/components/classic/}.

\bibitem{jmh}
OpenJDK.
\newblock \emph{Java Microbenchmark Harness (JMH)}, version 1.37.
\newblock \url{https://github.com/openjdk/jmh}.

\bibitem{little}
J.~D.~C. Little.
\newblock A proof for the queuing formula $L = \lambda W$.
\newblock \emph{Operations Research}, 9(3):383--387, 1961.

\end{thebibliography}

\appendix
\clearpage
\section{Reproducibility}
\label{app:runs}

Every run lives in \code{benchmarks/runs/<run_id>/} of the repository, and its \code{run_id} ends with the commit it
ran from. Each run keeps its \code{manifest.json} (commit and clean tree, exact command, host, Docker resources, JVM
flags, every \code{TARANTINO_*} variable, and replicas and threads per trial), \code{corpus.json}, \code{log.txt} with
the hypotheses printed before measuring, the logs of every service, the raw data and \code{summary.csv}. Every cited
run comes from a clean tree on the official machine. The campaign that produced them is
\run{20261105-215348-campaign-c1376c7} (\code{./scripts/bench/all.sh}); its \code{runs.txt} lists the 187 steps with
their exit codes. \run{20261106-123526-campaign-b4667d8} (\code{./scripts/bench/all.sh --only broker-latency}) reran the
one step that failed.

\begin{table}[!ht]
\centering
\footnotesize
\caption{Source of every table and figure of Section~\ref{sec:benchmarks} and of Table~\ref{tab:decision}.}
\label{tab:runs}
\begin{tabularx}{\textwidth}{@{}L{2.2cm}L{5.4cm}L{3.6cm}Y@{}}
\toprule
In this report & Run & Command & Raw data \\
\midrule
Tables~\ref{tab:ingest} and \ref{tab:ingest-workers}; crawler JMH & \run{20261105-215351-ingest-throughput-c1376c7} & \code{./scripts/bench/ingest-throughput.sh} & \code{ingest.csv}, \code{workers.csv}, \code{stats.csv}, \code{host.csv}, \code{jmh/crawler.csv} \\
Figure~\ref{fig:index} & \run{20261105-220915-index-throughput-c1376c7} & \code{./scripts/bench/index-throughput.sh --without-jmh} & \code{index.csv}, \code{stats.csv} \\
Indexer JMH & \run{20261103-204454-index-throughput-6925d4e} & \code{./scripts/bench/index-throughput.sh} & \code{jmh/indexer.json}, \code{jmh/indexer.csv} \\
Figure~\ref{fig:search}, Table~\ref{tab:search-latency}; query JMH & \run{20261104-140900-search-latency-8e1c1e8} & \code{./scripts/bench/search-latency.sh} & \code{load/<trial>.csv.gz}, \code{jmh/query.csv}, \code{explain/} \\
Tables~\ref{tab:decision}, \ref{tab:coordination}; Figure~\ref{fig:modes} & \run{20261106-012054-coordination-variants-c1376c7} & \code{./scripts/bench/coordination-variants.sh} & \code{coordination.csv}, \code{visibility.csv}, \code{stats.csv}, \code{variants.csv}, \code{pairs.csv} \\
Table~\ref{tab:resources} & \run{20261105-235758-resource-usage-c1376c7} & \code{./scripts/bench/resource-usage.sh} & \code{stats.csv}, \code{host.csv} \\
Table~\ref{tab:visibility} & \run{20261106-002208-publication-visibility-c1376c7} & \code{./scripts/bench/publication-visibility.sh} & \code{visibility.csv}, \code{coordination.csv} \\
Figure~\ref{fig:broker} & \run{20261106-123534-broker-latency-b4667d8} & \code{./scripts/bench/broker-latency.sh} & \code{broker.csv}, \code{load/} \\
Table~\ref{tab:recovery} & The 165 runs \code{...-recovery-<nn>[-<mode>]-r<N>-c1376c7}, listed in \code{sources.txt} of \run{20261106-073908-recovery-summary-c1376c7} & \code{TARANTINO_VARIANT=<mode> ./scripts/recovery/<nn>-<row>.sh <N>} & \code{checks.json}, \code{summary.csv} \\
Table~\ref{tab:delay} & \run{20261106-073911-recovery-delay-c1376c7} & \code{./scripts/bench/recovery-delay.sh <recovery runs>} & \code{recovery-delay.csv} \\
\bottomrule
\end{tabularx}
\end{table}

Each table and chart of this report takes its numbers from the \code{summary.csv}, \code{variants.csv},
\code{pairs.csv} or \code{workers.csv} of the run named in its caption. The rows they use are written into this
report's source, \code{QueryTarantino_Stage2_Report.tex}, at the root of the repository. The medians of
Table~\ref{tab:delay} come from \code{recovery-delay.csv}, over the books of every run of each row, and in row 12 apart
by their \code{status}. \code{bench ingest-workers} and \code{bench resources} added \code{workers.csv} and the CPU rows
of the ingestion run after the campaign, from the transition lines and the \code{stats.csv} the run already kept. To
summarize a run again, delete its \code{summary.csv} and run its \code{bench} commands on the run folder, for example:
\begin{lstlisting}
java -jar bench/target/bench-1.0-SNAPSHOT.jar coordination benchmarks/runs/<run_id>
\end{lstlisting}
\code{bench charts} redraws the charts of a run. The report itself builds from the root of the repository with:
\begin{lstlisting}
tectonic -X compile QueryTarantino_Stage2_Report.tex
\end{lstlisting}
\clearpage
\section{Endpoint Inventory}
\label{app:endpoints}

Table~\ref{tab:endpoints} lists every HTTP operation of \code{api/openapi.yaml} (\spec{6.1}). Besides the responses
listed, every endpoint answers \code{500} with Problem Details to an unexpected exception. Every \code{202} carries a
\code{Location} header to poll, and every \code{503} carries \code{Retry-After: 5}, or \code{1} on the internal routes.

\begin{table}[!ht]
\centering
\footnotesize
\caption{Every HTTP operation, with its provider, nginx port, consumer and responses.}
\label{tab:endpoints}
\begin{tabularx}{\textwidth}{@{}L{1.75cm}L{1.95cm}L{0.95cm}L{4.15cm}L{2.6cm}Y@{}}
\toprule
Inventory & Provider & nginx & Method and path & Consumer & Responses \\
\midrule
Front-end & \code{query} (7003) & 8080 & \code{GET /search} & Front-end & 200, 400, 503 \\
Front-end & \code{query} (7003) & 8080 & \code{GET /books/{book_id}} & Front-end & 200, 400, 404, 503 \\
Front-end & \code{query} (7003) & 8080 & \code{GET /search/filters} & Front-end & 200, 503 \\
Operational & \code{crawler} (7001) & 8081 & \code{POST /ingest/{book_id}} & Load and admin scripts & 202, 200, 400, 503 \\
Operational & \code{crawler} (7001) & 8081 & \code{GET /ingest/status/{book_id}} & Whoever got a 202; checks & 200, 400, 404, 503 \\
Operational & \code{crawler} (7001) & 8081 & \code{GET /ingest/list} & Checks & 200, 503 \\
Operational & \code{indexer} (7002) & 8081 & \code{POST /index/update/{book_id}} & Admin & 202, 200, 400, 404, 409, 503 \\
Operational & \code{indexer} (7002) & 8081 & \code{GET /index/status/{book_id}} & Whoever got a 202 & 200, 400, 404, 503 \\
Operational & \code{indexer} (7002) & 8081 & \code{POST /index/rebuild} & Admin & 202, 503 \\
Operational & \code{indexer} (7002) & 8081 & \code{GET /index/status} & Admin, benchmarks & 200, 503 \\
Operational & \code{control} (7004) & 8081 & \code{GET /control/status} & Admin, benchmarks & 200, 503 \\
Health & Every service & -- & \code{GET /health} & Compose healthchecks & 200, 503 \\
Internal & \code{crawler} (7001) & 8082 & \code{POST /internal/download} & \code{control}, mode \mode{orchestrator} & 202, 400, 503 \\
Internal & \code{indexer} (7002) & 8082 & \code{POST /internal/index} & \code{control}, mode \mode{orchestrator} & 202, 400, 503 \\
Internal & Each \code{query} & -- & \code{POST /internal/snapshot/refresh} & \code{control}, mode \mode{orchestrator} & 202 \\
\bottomrule
\end{tabularx}
\end{table}

\begin{table}[!ht]
\centering
\small
\caption{Answer of each state-changing request, by the book's state (\spec{6.2}). Every response body is the book's
state: \code{book_id}, \code{status}, \code{attempts}, \code{last_error} and \code{updated_at}.}
\label{tab:state-answers}
\begin{tabularx}{\textwidth}{@{}L{5.6cm}Y@{}}
\toprule
Request and book state & Response \\
\midrule
\code{POST /ingest}, no document & ``Create the state'' and enqueue: \code{202} with \code{QUEUED} \\
\code{POST /ingest}, \code{FAILED} & ``Retry a \code{FAILED} book'': \code{202} with \code{QUEUED} and \code{attempts} back to 0 \\
\code{POST /ingest}, \code{QUEUED} or \code{DOWNLOADING} & \code{202} with its state, without enqueuing it again \\
\code{POST /ingest}, \code{DOWNLOADED}, \code{INDEXING} or \code{PUBLISHED} & \code{200} with its state \\
\code{POST /index/update}, no document & \code{404 BOOK_NOT_FOUND} \\
\code{POST /index/update}, \code{QUEUED}, \code{DOWNLOADING} or \code{FAILED} & \code{409 NOT_DOWNLOADED}; a \code{FAILED} book goes forward with \code{POST /ingest} \\
\code{POST /index/update}, \code{DOWNLOADED} & \code{202} with its state, and a \code{BookTask} to \code{books-to-index}; a duplicate does no harm \\
\code{POST /index/update}, \code{INDEXING} & \code{202} with its state, without enqueuing \\
\code{POST /index/update}, \code{PUBLISHED} & \code{200} with its state \\
\bottomrule
\end{tabularx}
\end{table}
\clearpage
\section{Failure Matrix}
\label{app:failures}

Each row of Table~\ref{tab:failures} has its own script, \code{scripts/recovery/<nn>-<row>.sh <N>}. The script brings
the stack up, causes the failure, restarts what went down, and checks three things: every book ends
\code{PUBLISHED}, \code{/search} finds it exactly once, and no book is in two published batches of the current
generation (\spec{8}). A crash point stops the service with \code{Runtime.halt(1)}, without shutdown hooks, the first
time it reaches the named point.

\begin{table}[!ht]
\centering
\footnotesize
\caption{What each failure leaves behind, who notices and how the system recovers, in mode \mode{hazelcast}. In the
other modes, rows 04, 05 and 12 recover by redelivery or a new dispatch, or do not apply (\spec{11.5}).}
\label{tab:failures}
\begin{tabularx}{\textwidth}{@{}L{0.65cm}L{2.75cm}L{2.85cm}L{2.2cm}YL{2.65cm}@{}}
\toprule
Row & What goes down & What is left & Who notices & How it recovers & How it is caused \\
\midrule
01 & \code{crawler} while downloading & A \code{DOWNLOADING} book with a lease & \code{control}, when it expires & Back to \code{QUEUED} and to the queue & \code{after-claim} \\
02 & \code{crawler} after the datalake write & The book in the datalake, still \code{DOWNLOADING} & \code{control}, when it expires & Back to \code{QUEUED}; the next claim finds it in the datalake & \code{after-datalake-write} \\
03 & \code{crawler} after \code{DOWNLOADED}, before enqueuing & A \code{DOWNLOADED} book without a message & \code{control}: empty queue, old \code{enqueued_at} & Enqueued again & \code{after-downloaded} \\
04, 05 & \code{crawler} or \code{indexer} after taking messages & Books without a message: an \code{IQueue} has no acknowledgement & \code{control}: empty queue, old \code{enqueued_at} & Enqueued again & \code{after-dequeue} \\
06 & \code{indexer} while writing a batch & A \code{WRITING} batch, \code{INDEXING} books, half-written postings & \code{control}, when the batch expires & \code{ABANDONED}; books back to the queue; postings deleted later & \code{during-batch-write} \\
07 & \code{indexer} after publishing, before marking & A \code{PUBLISHED} batch with \code{INDEXING} books & \code{control} & Completes the publication & \code{after-publish} \\
08 & \code{indexer} after marking, before notifying & A published batch without \code{BatchPublished} & \code{query}, on its refresh & Seen within 5\,s & \code{after-mark} \\
09 & \code{crawler} stopped longer than its lease & A zombie worker that may still write & \code{control}, when the lease expires & Requeued; the zombie's writes fail on \code{claim_id}, and its rename only replaces a complete copy & \code{docker pause} for 10\,s \\
10 & \code{indexer} stopped longer than its lease & A zombie indexer that keeps writing & Its own renewal fails & It stops writing; its publication fails, and \code{control} deletes its postings, also late ones & \code{docker pause} for 10\,s \\
11 & \code{control} & At most books in a new state without their message & Its reconciliation, once back & Reconciles from MongoDB and requeues & \code{docker kill} while seeding \\
12 & Hazelcast & Empty queues & \code{control}: empty queue, old \code{enqueued_at} & Requeues every waiting book & \code{docker restart} \\
13 & MongoDB & No half-done transition: each one is a single write & Healthchecks; services answer \code{503} & Everything continues from the last state & \code{docker stop}, 10\,s \\
14 & One \code{query} replica & Nothing: it does not write & nginx & The other replicas serve; it rebuilds its snapshot at startup & \code{docker kill} \\
\bottomrule
\end{tabularx}
\end{table}
\clearpage
\section{Configuration}
\label{app:config}

Every service reads its configuration from environment variables through \code{TarantinoConfig}, which validates
them at startup and stops the service naming any invalid one (\spec{4.7}). Table~\ref{tab:variables} lists the
variables that set ports, addresses, data paths, replicas, threads and timing. Compose passes the same timing values to
every service, and \code{scripts/recovery/recovery.env} shortens them for the recovery tests while keeping the
invariants: heartbeat $\leq$ lease / 3, garbage collection and orphan age $\geq$ 10 leases, and requeue $>$ batch
timeout.

\begin{table}[!ht]
\centering
\footnotesize
\caption{Main configuration variables, their defaults and the recovery-test values.}
\label{tab:variables}
\begin{tabularx}{\textwidth}{@{}L{5.55cm}L{2.75cm}L{1.3cm}Y@{}}
\toprule
Variable & Default & Recovery & Purpose \\
\midrule
\code{TARANTINO_PORT} & 7001--7004 & & HTTP port of \code{crawler}, \code{indexer}, \code{query}, \code{control} \\
\code{TARANTINO_HTTP_THREADS} & Javalin's, 8--250 & & Maximum request threads \\
\code{TARANTINO_MONGO_URI}, \code{TARANTINO_HAZELCAST} & \code{localhost} & & MongoDB and Hazelcast addresses; Compose points them to its services \\
\code{TARANTINO_DATALAKE}, \code{TARANTINO_WORKLOAD}, \code{TARANTINO_MIRROR} & \code{datalake}, \code{workload}, none & & Datalake root, workload folder and local Gutenberg mirror \\
\code{TARANTINO_CANDIDATES} & \code{sample_ids.txt} & & Books that \code{control} seeds, or \code{mirror} \\
\code{TARANTINO_PARALLEL_DOWNLOADS} & 8 & & Download workers per \code{crawler} replica \\
\code{TARANTINO_MAX_ATTEMPTS} & 3 & & Downloads of a book before \code{FAILED} \\
\code{TARANTINO_INDEX_BATCH} & 100 & 5 & Books per batch \\
\code{TARANTINO_BATCH_TIMEOUT_SECONDS} & 30 & 5 & Longest wait to close a batch \\
\code{TARANTINO_LEASE_SECONDS} & 60 & 6 & Duration of a book or batch lease \\
\code{TARANTINO_HEARTBEAT_SECONDS} & 20 & 2 & How often a lease is renewed \\
\code{TARANTINO_RECONCILE_SECONDS} & 30 & 3 & How often \code{control} reconciles \\
\code{TARANTINO_REQUEUE_SECONDS} & 120 & 10 & Age after which a lost message is requeued \\
\code{TARANTINO_GC_SECONDS} & 600 & 60 & Wait before deleting abandoned postings and old generations \\
\code{TARANTINO_SNAPSHOT_SECONDS} & 5 & 2 & Periodic refresh of the \code{query} snapshot \\
\code{TARANTINO_CORS_ORIGIN} & \code{*} & & Origin that \code{query} allows \\
\code{TARANTINO_COORDINATION} & \code{hazelcast} & & Coordination mode \\
\code{TARANTINO_SNAPSHOT_ON_PUBLISH} & \code{true} & & Topic listener in \code{query} \\
\code{TARANTINO_ACTIVEMQ_PERSISTENT} & \code{false} & & Persistent delivery and acknowledgement after the job \\
\code{TARANTINO_POLL_MILLIS} & 200 & & Sleep after an empty poll, mode \mode{polling} \\
\code{TARANTINO_ORCHESTRATOR_TICK_MILLIS} & 200 & & Pause between ticks, mode \mode{orchestrator} \\
\code{JAVA_TOOL_OPTIONS} & \code{-Xmx512m}, exit when out of memory & & Fixed heap; a JVM out of memory exits instead of losing a thread \\
\bottomrule
\end{tabularx}
\end{table}

\end{document}
